\chapter{评估与安全：如何知道模型足够好？}
\label{ch:eval}
%% ============================================================

\section{评估的困境}

评估 LLM 是一个根本性的难题。
传统的 NLP 评估（给定答案，检查是否正确）
对于开放式生成任务完全不适用。
一个问题可以有无数种正确的回答方式。

\subsection{Goodhart 定律}

\begin{warnbox}[Goodhart 定律的警告]
「当一个度量成为目标时，它就不再是好的度量。」
已有证据表明，一些模型在特定基准测试上被过度优化，
在 MMLU 上得分很高，但在实际使用中表现平平。
Chatbot Arena 的人类盲评被认为是最可靠的评估方式，
因为它直接衡量用户满意度，
而且每天更新的问题使得「针对性优化」变得困难。
\end{warnbox}

Goodhart 定律在 LLM 评估中有多种表现形式：

\textbf{直接过拟合}：
当模型开发者将 MMLU 得分作为关键指标反复迭代时，
他们可能有意或无意地调整训练数据、训练策略和超参数
来最大化 MMLU 分数。
结果是模型学会了「做 MMLU 选择题」的模式，
而不是真正掌握了 57 个学科的知识。

\textbf{格式 hacking}：
某些基准测试对输出格式敏感。
例如，MMLU 期望模型输出 A/B/C/D 中的一个字母。
一些模型被特别训练成在遇到选择题时
以极高的概率直接输出字母，
而不是先「思考」再回答。
这提高了基准分数，但可能降低了模型在
需要深思熟虑的真实场景中的表现。

\textbf{奖励模型 hacking}：
在使用 RLHF 时，模型可能学会了讨好奖励模型
（生成奖励模型喜欢的风格，如冗长、过度谨慎、
频繁使用“I appreciate your question” 这样的客套话），
而不是真正提高回答质量。
这是 Goodhart 定律在训练阶段的体现。

\subsection{数据污染}
\label{subsec:data-contamination}

\textbf{数据污染}（data contamination）
是基准测试面临的另一个严峻挑战：
如果测试题目出现在了模型的训练数据中，
模型可能只是在「背答案」而非真正推理。

随着训练数据规模从数十亿扩大到数万亿 token，
几乎不可能保证训练数据中完全不包含任何基准测试的内容，
因为这些测试题本身就发布在互联网上。

常用的检测和缓解方法包括：
\begin{itemize}[noitemsep]
  \item \textbf{N-gram overlap detection}：
        检查训练数据和测试题之间的 n-gram 重叠率。
        超过阈值的样本被标记为可能泄露。
  \item \textbf{Canary strings}：
        在测试集中插入独特的标记字符串，
        然后检查模型是否能记忆这些字符串。
  \item \textbf{Temporal splits}：
        使用模型训练截止日期\textbf{之后}发布的数据构建测试集。
        LiveBench 就采用这种策略，定期更换测试题。
  \item \textbf{Private test sets}：
        不公开测试题的全部内容。
        SWE-bench 的部分测试用例就是私有的。
\end{itemize}

\subsubsection{深入数据污染检测}

近年来，数据污染检测已发展为一个独立的研究方向。
Ravaut et al.\ (2024)~\cite{ravaut2024contamination}
对现有方法进行了系统综述，
将检测方法分为三个层面：

\textbf{第一层：基于训练数据的检测}。
如果能访问模型的训练数据（通常仅限开源模型），
可以直接进行文本匹配。
方法包括精确字符串匹配、
N-gram 重叠率计算（通常用 8-gram 或 13-gram），
以及语义相似度检索。
GPT-4 的技术报告就使用了
基于子串匹配的污染检测，
但由于训练数据本身不公开，
外部研究者无法独立验证。

\textbf{第二层：基于模型行为的检测}。
即使无法访问训练数据——
也可以通过分析模型行为推断是否存在污染：
\begin{itemize}[noitemsep]
  \item \textbf{成员推断攻击}（Membership Inference Attack, MIA）：
        利用模型对训练样本和非训练样本的不同行为
        （如 perplexity 差异）来判断某个样本是否在训练集中。
        Shi et al.\ (2024) 提出的 Min-K\% Prob 方法~\cite{shi2024detecting}
        通过检查一段文本中概率最低的 K\% token
        的平均对数概率来判断：
        如果这些「最不可能」的 token 的概率仍然很高，
        说明模型很可能在训练中见过这段文本。
  \item \textbf{补全测试}（Completion Test）：
        给模型一道测试题的前半部分，
        检查它能否逐字补全后半部分。
        如果模型能精确复现原文的表述
        （而非给出语义等价但措辞不同的回答），
        强烈暗示该题目出现在训练数据中。
  \item \textbf{扰动一致性测试}（Perturbation Test）：
        对基准测试题做微小的\textbf{语义保持}修改
        （改变数字、替换名称、调整选项顺序），
        如果模型在原题上准确率很高
        但在扰动版本上大幅下降，
        说明它可能只是记住了原题的答案。
        Golchin \& Surdeanu (2024)~\cite{golchin2024time}
        系统地研究了这种方法。
\end{itemize}

\textbf{第三层：基于基准设计的预防}。
从源头上减少污染的可能性：
\begin{itemize}[noitemsep]
  \item \textbf{动态基准}：
        LiveBench 和 LiveCodeBench 定期更新测试题，
        使用模型训练截止日期之后的数据，
        从根本上消除了数据泄露的可能性。
  \item \textbf{私有测试集}：
        保留一部分测试题不公开。
        只接受 API 提交，服务端验证结果。
  \item \textbf{程序化生成}：
        使用参数化模板动态生成测试题。
        每次评估使用不同的具体实例，
        但考察相同的能力。
        FrontierMath 的部分题目就采用了这种策略。
\end{itemize}

\subsubsection{污染的产业化与系统性风险}

数据污染问题在 2025--2026 年呈现出新的形态，
从偶然的「无意泄露」演变为系统性的产业风险。
这种转变的根本原因在于：
当基准分数直接影响融资额、市场估值和客户决策时，
「优化基准」的激励变得不可忽视。

\textbf{合成数据的反污染陷阱}尤其值得关注。
随着合成数据在训练中的比例不断上升~\cite{liu2024synthetic}，
一种新的污染路径出现了：
如果用强模型（如 GPT-4o）生成合成训练数据，
而这些强模型本身已经在基准测试上被评估过且记住了部分题目，
那么生成的合成数据可能\textbf{间接包含基准测试的信息}。
这种「二阶污染」极难检测，
因为合成文本与原始测试题在表面上没有任何 n-gram 重叠，
但在语义层面可能编码了相同的知识模式。
2025 年的一项研究~\cite{mehta2025contamination}
在多个基准上验证了这种效应的存在：
使用 GPT-4o 生成的数学练习题
与 MATH 基准的部分题目在解题策略上
呈现出统计可辨的相似性。

\textbf{跨基准的污染传导}是另一个被忽视的问题。
许多基准测试共享来源（如 Wikipedia、学术论文、教科书），
一个模型在 MMLU 上被检测为「无污染」
不代表其训练数据中不包含与 MMLU 题目高度相关的材料。
一个数学定理的证明过程可能同时出现在
MATH、AIME 和 GPQA 的多道题目的参考解法中。
训练数据包含了这个证明，
模型就可能在多个基准上同时获得不公平的优势。

\begin{warnbox}[SWE-bench Verified 的污染警示]
2026 年 2 月，OpenAI 发布报告指出
SWE-bench Verified 已出现严重的数据污染问题~\cite{openai2026swebenchpro}。
由于该基准的测试用例来自公开的 GitHub issue，
许多模型的训练数据中包含了这些 issue 的讨论和解决方案。
OpenAI 建议社区转向 SWE-bench Pro，
一个使用私有测试用例和更严格防泄露措施的新版本。
这一事件凸显了\textbf{任何公开的静态基准
都有被污染的风险}，
进一步强化了动态基准和私有测试集的重要性。
\end{warnbox}

\subsection{基准饱和}

LLM 的进步速度让基准测试快速\textbf{饱和}：
设计来区分好坏的测试，
在最新的模型面前失去了区分度。

GSM8K（8500 道小学数学应用题）在 2023 年初
还被认为是衡量数学推理能力的有力基准。
到 2025 年底，前沿模型在 GSM8K 上的准确率普遍超过 95\%，
GSM8K 已经无法区分不同模型的数学能力。
社区因此推出了更难的替代品：
MATH（竞赛级数学）$\to$ AIME（数学奥林匹克）$\to$
GPQA（博士级科学问题）$\to$ FrontierMath（研究级数学）。
每一代更难的基准测试的生命周期都在缩短。

这导致了一场\textbf{基准军备竞赛}：
研究者不断设计更难的测试，模型不断攻克它们。
这本身说明了 LLM 能力的快速进步，
但也给评估方法论带来了挑战：
我们需要的不仅是更难的题目，
还需要更好的评估\textbf{范式}。

\subsubsection{传统基准的饱和（2026 年现状）}

2026 年初，基准饱和问题变得更加严峻。
多个前沿模型在 \textbf{AIME 2025} 上的准确率超过 96\%：
Doubao Seed 2.0 Pro 达到 98.3\%，
Step 3.5 Flash 达到 97.3\%，
Kimi K2.5 达到 96.1\%，
Qwen3-Max-Thinking 在部分配置下甚至报告 100\%。
当多个模型都能做对几乎所有题目时，
基准就失去了区分度。

社区正在转向更难的替代品作为新的区分器：
\textbf{HMMT}（Harvard-MIT Mathematics Tournament）
和 \textbf{GPQA Diamond} 成为区分前沿模型的新标准。
而 \textbf{HLE}（Humanity's Last Exam）
代表了评估的最远边界。
GLM-4.7 在 HLE 上仅达到 42.8\%，
说明即使是最强的 frontier 模型
距离全面的人类专家水平仍然相当遥远。

\subsubsection{Humanity's Last Exam：评估的终极试验}

\textbf{Humanity's Last Exam}（HLE）~\cite{phan2025hle}
是由 Scale AI 和 Center for AI Safety 于 2025 年联合发起的
超高难度基准测试，其设计理念是：
创造一组\textbf{当前最强 AI 系统也无法通过的考试}。

HLE 的独特之处在于它的题目来源：
约 3,000 道题目由全球各领域的顶尖专家贡献，
涵盖数学、物理、生物、哲学、法律、艺术史等
超过 100 个专业领域。
每道题目都经过严格筛选：
如果任何一个参与评估的 frontier 模型能够正确回答，
该题目就被排除。
这种「对抗性筛选」确保了 HLE 中的每道题目
都位于当前 AI 能力的\textbf{真正边界之外}。

HLE 的题目质量由一个
多层次的同行评审流程保障：
每道题目至少经过 3 位独立专家的审核，
确保题意明确、答案唯一且可验证、
难度确实超越了当前 AI 系统的能力。
此外，HLE 引入了一个创新的
\textbf{跨学科验证}机制：
每道题不仅由同领域专家审核，
还由其他领域的专家
从「外行视角」检验题目是否
可以通过通用推理能力解决
（如果可以，则该题目不够难，
因为 AI 系统的通用推理能力可能足以应对）。

截至 2026 年 3 月，HLE 的状态值得深思：
\begin{itemize}[noitemsep]
  \item 最强的推理模型在 HLE 上的准确率约为 35--43\%。
  \item 标准模型（非推理模式）的准确率低于 20\%。
  \item 人类专家（在各自领域内）的准确率约为 70--80\%。
  \item 题目中约 15\% 被归类为「开放性问题」，
        即使是人类专家也没有确定答案。
  \item 推理模型（如 o3）相比标准模型（如 GPT-4o）
        有约 10--15 个百分点的提升，
        但远未达到人类专家水平。
\end{itemize}

HLE 的存在回答了一个关键问题：
\textbf{AI 离「全面超越人类专家」还有多远？}
答案是，在需要深度专业知识的领域，
差距仍然显著。
但 HLE 也面临自身的局限性：
如果 AI 能力继续按当前速度提升，
HLE 可能在 2--3 年内也面临饱和风险，
届时社区需要设计更具挑战性的评估方式，
或者承认「静态考试」作为评估范式的根本局限。

\subsubsection{OpenRouter 匿名模型评测}

一种全新的评估范式正在 OpenRouter 平台上自发形成：
\textbf{匿名模型投放}（stealth model deployment）。
模型提供者以匿名身份将新模型发布到 OpenRouter，
用户在不知道模型身份的情况下进行测试和反馈，
完全消除了品牌偏见对评估结果的影响。

这一实践始于 2025 年底的 \textbf{Quasar Alpha}——
随后出现了 Optimus Alpha、Aurora Alpha、Pony Alpha 等匿名模型，
每一个都引发了社区的激烈讨论和身份猜测。
2026 年 3 月 11 日发布的 \textbf{Hunter Alpha}
号称具有「1T 参数 + 1M 上下文」的规格，
其参数规格与预期中的 DeepSeek V4 高度吻合。
同期发布的 \textbf{Healer Alpha} 被描述为
“frontier omni-modal model”。

匿名模型评测的价值在于：
它提供了一个\textbf{品牌无关}的能力评估环境。
用户对模型的评价完全基于输出质量，
不受「这是 GPT 还是 Claude」的先验预期影响。
这种评估方式可以看作是
Chatbot Arena 盲评理念的一种自然延伸。

\subsubsection{Arena-Lite：基于锦标赛的评估}

\textbf{Arena-Lite}（EMNLP 2025）提出了
一种基于\textbf{锦标赛制}的直接比较评估方法。
与传统评估依赖基线模型输出不同，
Arena-Lite 让模型之间进行\textbf{直接对决}——
消除了对参考输出的依赖。
这解决了 AlpacaEval 等基准
受限于基线模型质量的问题：
当基线模型本身不够好时，
评估结果的可信度也随之降低。

截至 2026 年 3 月，各主要评估平台上
排名最前的模型包括：
Claude Opus 4.6、GPT-5.4、DeepSeek V3.2、
Gemini 3.1 Pro、Qwen 3.5。

\section{基准测试全景}

主流基准测试覆盖不同能力维度。
为了系统化地理解这一格局，
图~\ref{fig:benchmark-taxonomy} 以层次结构呈现了
当前主要基准测试的分类体系。

%% --- TikZ: Benchmark Landscape Taxonomy ---
\begin{figure}[htbp]
\centering
\begin{tikzpicture}[
  >={Stealth[length=2.5pt,width=2.5pt]},
  every node/.style={font=\small},
  root/.style={rectangle, rounded corners=4pt, draw=figBlue!60,
    fill=figBlue!8, minimum height=9mm, align=center,
    font=\small\bfseries, inner sep=6pt},
  cat/.style={rectangle, rounded corners=3pt, draw=#1!50,
    fill=#1!8, minimum height=8mm, text width=20mm,
    align=center, font=\scriptsize\bfseries, inner sep=4pt},
  pill/.style={rectangle, rounded corners=6pt, draw=#1!40,
    fill=#1!6, minimum height=5.5mm, inner sep=3pt,
    font=\tiny, align=center},
  conn/.style={draw=figGray!40, line width=0.5pt},
]

% Root node
\node[root] (root) at (0, 0) {LLM 基准测试全景};

% Category nodes - spread across width
\node[cat=figRed]    (know)   at (-5.6, -1.6) {知识};
\node[cat=figOrange] (reason) at (-2.8, -1.6) {推理};
\node[cat=figGreen]  (code)   at (0, -1.6)    {代码};
\node[cat=figGreen]  (safe)   at (2.8, -1.6)  {安全};
\node[cat=figBlue]   (agent)  at (5.6, -1.6)  {Agent};

% Connections root -> categories
\foreach \cat in {know, reason, code, safe, agent} {
  \draw[conn] (root.south) -- (\cat.north);
}

% --- Leaf benchmarks as pills ---
% Knowledge (saturated - figRed)
\node[pill=figRed] (mmlu)    at (-6.4, -3.0) {MMLU};
\node[pill=figRed] (mmlup)   at (-6.4, -3.7) {MMLU-Pro};
\node[pill=figRed] (arc)     at (-4.8, -3.0) {ARC};
\node[pill=figRed] (gpqa)    at (-4.8, -3.7) {GPQA};
\draw[conn] (know.south) -- (mmlu.north);
\draw[conn] (know.south) -- (mmlup.north);
\draw[conn] (know.south) -- (arc.north);
\draw[conn] (know.south) -- (gpqa.north);

% Reasoning (near-saturated - figOrange)
\node[pill=figOrange] (gsm)   at (-3.6, -3.0) {GSM8K};
\node[pill=figOrange] (math)  at (-3.6, -3.7) {MATH};
\node[pill=figOrange] (aime)  at (-2.0, -3.0) {AIME};
\node[pill=figOrange] (fmath) at (-2.0, -3.7) {FrontierMath};
\draw[conn] (reason.south) -- (gsm.north);
\draw[conn] (reason.south) -- (math.north);
\draw[conn] (reason.south) -- (aime.north);
\draw[conn] (reason.south) -- (fmath.north);

% Code (discriminative - figGreen)
\node[pill=figGreen] (heval) at (-0.8, -3.0) {HumanEval};
\node[pill=figGreen] (mbpp)  at (-0.8, -3.7) {MBPP};
\node[pill=figGreen] (swe)   at (0.8, -3.0)  {SWE-bench};
\node[pill=figGreen] (lcb)   at (0.8, -3.7)  {LiveCodeBench};
\draw[conn] (code.south) -- (heval.north);
\draw[conn] (code.south) -- (mbpp.north);
\draw[conn] (code.south) -- (swe.north);
\draw[conn] (code.south) -- (lcb.north);

% Safety (evolving - figGreen)
\node[pill=figGreen] (tqa)  at (2.0, -3.0) {TruthfulQA};
\node[pill=figGreen] (bbq)  at (2.0, -3.7) {BBQ};
\node[pill=figGreen] (harm) at (3.6, -3.0) {HarmBench};
\node[pill=figGreen] (advb) at (3.6, -3.7) {AdvBench};
\draw[conn] (safe.south) -- (tqa.north);
\draw[conn] (safe.south) -- (bbq.north);
\draw[conn] (safe.south) -- (harm.north);
\draw[conn] (safe.south) -- (advb.north);

% Agent (emerging - figBlue)
\node[pill=figBlue] (weba) at (4.8, -3.0) {WebArena};
\node[pill=figBlue] (osw)  at (4.8, -3.7) {OSWorld};
\node[pill=figBlue] (bfcl) at (6.4, -3.0) {BFCL};
\node[pill=figBlue] (tbch) at (6.4, -3.7) {ToolBench};
\draw[conn] (agent.south) -- (weba.north);
\draw[conn] (agent.south) -- (osw.north);
\draw[conn] (agent.south) -- (bfcl.north);
\draw[conn] (agent.south) -- (tbch.north);

% --- Saturation status labels (white fill to avoid overlap with connection lines) ---
\node[font=\tiny, figRed, anchor=north, fill=white, inner sep=1pt, rounded corners=0.5pt]
  at (know.south)   [yshift=-0.5mm] {饱和};
\node[font=\tiny, figOrange, anchor=north, fill=white, inner sep=1pt, rounded corners=0.5pt]
  at (reason.south) [yshift=-0.5mm] {接近饱和};
\node[font=\tiny, figGreen, anchor=north, fill=white, inner sep=1pt, rounded corners=0.5pt]
  at (code.south)   [yshift=-0.5mm] {区分度高};
\node[font=\tiny, figGreen, anchor=north, fill=white, inner sep=1pt, rounded corners=0.5pt]
  at (safe.south)   [yshift=-0.5mm] {持续演进};
\node[font=\tiny, figBlue, anchor=north, fill=white, inner sep=1pt, rounded corners=0.5pt]
  at (agent.south)  [yshift=-0.5mm] {新兴};

% --- Legend ---
\node[anchor=north west, inner sep=0pt] at (-7.0, -4.6) {%
  \begin{tikzpicture}[every node/.style={font=\tiny, inner sep=0pt}]
    \fill[figRed!8]    (0, 0)     rectangle (0.25, 0.2);
    \draw[figRed!40]   (0, 0)     rectangle (0.25, 0.2);
    \node[right=2pt] at (0.25, 0.1) {饱和};
    %
    \fill[figOrange!8] (1.3, 0)   rectangle (1.55, 0.2);
    \draw[figOrange!40](1.3, 0)   rectangle (1.55, 0.2);
    \node[right=2pt] at (1.55, 0.1) {接近饱和};
    %
    \fill[figGreen!8]  (3.0, 0)   rectangle (3.25, 0.2);
    \draw[figGreen!40] (3.0, 0)   rectangle (3.25, 0.2);
    \node[right=2pt] at (3.25, 0.1) {区分度高};
    %
    \fill[figBlue!8]   (4.8, 0)   rectangle (5.05, 0.2);
    \draw[figBlue!40]  (4.8, 0)   rectangle (5.05, 0.2);
    \node[right=2pt] at (5.05, 0.1) {新兴};
  \end{tikzpicture}%
};

\end{tikzpicture}
\caption{LLM 基准测试全景分类。
颜色标注反映了各领域基准测试在 2026 年初的饱和状态：
红色表示已饱和、橙色接近饱和、绿色仍有区分度、蓝色为新兴领域。}
\label{fig:benchmark-taxonomy}
\end{figure}

\subsection{知识与推理}
\begin{itemize}[noitemsep]
  \item \textbf{MMLU}（57 学科选择题）~\cite{hendrycks2021mmlu}：
        从高中到专业级别的知识测试。
        涵盖人文、社会科学、STEM 等。
        2023 年的标杆，2025 年已接近饱和（前沿模型 $>$90\%）。
  \item \textbf{MMLU-Pro}~\cite{wang2024mmlupro}：MMLU 的加强版。
        10 个选项（而非 4 个），更难的题目，更少的格式泄露。
        更好地区分当前前沿模型的知识深度。
  \item \textbf{ARC}（AI2 Reasoning Challenge）：
        科学推理题。ARC-Challenge 子集需要多步推理。
  \item \textbf{GPQA}（Graduate-Level Google-Proof QA）~\cite{rein2024gpqa}：
        博士级别的物理、化学、生物问题。
        “Google-Proof” 意味着即使允许搜索也不容易回答——
        需要深度专业知识和推理。
  \item \textbf{Big-Bench Hard}：
        从 BIG-Bench 的 204 个任务中筛选出
        语言模型表现最差的 23 个任务。
        专注于需要多步推理的困难问题。
\end{itemize}

\subsubsection{MMLU-Pro：超越 MMLU 的下一代知识基准}

MMLU 在 2023--2024 年是 LLM 评估的核心指标，
但随着前沿模型准确率逼近 90\% 甚至更高，
其区分度急剧下降。
Wang et al.\ (2024) 在 NeurIPS 2024 上提出了 MMLU-Pro~\cite{wang2024mmlupro}，
通过三个关键改进重新拉开了模型之间的差距：

\begin{enumerate}[noitemsep]
  \item \textbf{10 选项制}：
        从 MMLU 的 4 个选项增加到 10 个。
        这将随机猜测的基线从 25\% 降低到 10\%，
        大幅减少了「蒙对」的概率。
        更重要的是，
        10 个选项中的干扰项（distractors）经过精心设计，
        反映了常见的错误推理路径，
        使得需要真正理解才能排除。
  \item \textbf{更高的推理需求}：
        MMLU-Pro 的题目需要更多的多步推理。
        实验表明，使用 Chain-of-Thought（CoT）提示
        可以将 MMLU-Pro 的准确率提升 10--15 个百分点，
        而在原始 MMLU 上 CoT 的提升幅度不到 2\%。
        这说明 MMLU-Pro 的题目确实在考察推理能力，
        而非简单的知识记忆。
  \item \textbf{更强的鲁棒性}：
        在 MMLU 上，选项顺序的变化可导致模型准确率波动 4--5\%；
        在 MMLU-Pro 上，这种波动降低到 2\% 以内。
        10 选项制降低了格式偏差（format bias）的影响。
\end{enumerate}

截至 2026 年初，前沿模型在 MMLU-Pro 上的准确率
约为 85--90\%（对比 MMLU 上的 90--95\%），
仍有足够的区分度。

\subsubsection{GPQA Diamond：博士级的「Google-Proof」评估}

GPQA~\cite{rein2024gpqa} 的设计理念是创造
即使允许使用搜索引擎也难以回答的问题。
题目由各领域的博士专家撰写并交叉验证：
\textbf{领域专家}的准确率约为 65\%，
而\textbf{拥有 Google 搜索权限的非专家}准确率仅约 34\%。
这个巨大的差距证明了题目确实需要深度专业知识。

GPQA Diamond 是 GPQA 中最具挑战性的 198 道题子集，
涵盖物理、化学和生物三个学科。
截至 2026 年初，
推理模型（如 o3、R1）在 GPQA Diamond 上的准确率
已达到 85--93\%，
\textbf{首次在某些领域超越了人类博士专家}。
这一方面说明了推理模型的巨大进步，
另一方面也意味着 GPQA 可能在未来 1--2 年内饱和。

\subsection{数学}
\begin{itemize}[noitemsep]
  \item \textbf{GSM8K}（8500 道小学应用题）：
        曾经的标准数学基准。需要 2--8 步的基本数学运算。
        2025 年已接近饱和，大多数前沿模型准确率 $>$95\%。
  \item \textbf{MATH}（12,500 道竞赛级数学）：
        覆盖代数、几何、数论等 7 个类别，
        难度从 1 到 5。Level 5 仍有挑战性。
  \item \textbf{AIME}（American Invitational Mathematics Examination）：
        美国数学邀请赛真题。整数答案，需要创造性数学推理。
        推理模型（R1, o1）在此基准上表现显著优于标准模型。
  \item \textbf{FrontierMath}~\cite{glazer2024frontiermath}：
        由 Epoch AI 与专业数学家合作设计的研究级数学基准。
        详见下文单独展开。
\end{itemize}

\subsubsection{FrontierMath：AI 数学能力的终极测试}

FrontierMath~\cite{glazer2024frontiermath}
由 Epoch AI 于 2024 年 11 月发布，
代表了目前最高难度的 AI 数学基准。
它与之前的数学基准有根本性的不同：

\begin{itemize}[noitemsep]
  \item \textbf{原创的未发表题目}：
        所有 350 道题目都是由数学家专门为该基准\textbf{原创}的，
        从未在任何地方发表过，
        从根本上杜绝了数据泄露。
  \item \textbf{研究级难度}：
        题目覆盖数论、实分析、代数几何、范畴论等
        现代数学的主要分支。
        解决一道典型的 FrontierMath 题目
        需要相关领域的数学研究者花费\textbf{数小时到数天}。
  \item \textbf{四级难度体系}：
        Tier 1--3 覆盖本科到博士后水平，
        Tier 4 是研究级数学。
        此外还有一组\textbf{开放问题}——
        即至今未被人类数学家解决的问题。
  \item \textbf{自动验证}：
        答案是精确的数值或数学对象，
        可以通过程序自动验证，无需主观判断。
\end{itemize}

\begin{keybox}[FrontierMath 的标志性意义]
当 FrontierMath 于 2024 年发布时，
所有前沿模型的通过率不到 2\%。
著名数学家 Terence Tao 评论道：
“这些题目极具挑战性——
我预计在数年内 AI 系统还无法解决其中的大部分。”
然而到 2025 年底，
推理模型（如 o3）在 Tier 1--2 上的通过率已超过 25\%，
在 Tier 4 上仍然低于 5\%。
FrontierMath 可能是目前唯一一个
在可预见的未来不会饱和的 LLM 基准。
\end{keybox}

\subsection{代码}
\begin{itemize}[noitemsep]
  \item \textbf{HumanEval}（164 道函数级编程题）：
        给出函数签名和 docstring，模型生成函数体。
        用单元测试验证正确性。简单但有用的基线。
  \item \textbf{MBPP}（Mostly Basic Python Problems）：
        974 道入门级 Python 编程题。比 HumanEval 更多样但也更简单。
  \item \textbf{SWE-bench}（真实 GitHub issue 修复）~\cite{jimenez2024swebench}：
        给模型一个真实的 GitHub issue 和整个代码仓库，
        要求生成修复 patch。
        这是目前最接近\textbf{真实软件工程}场景的基准，
        需要理解上千行代码、定位 bug、生成正确的修复。
        前沿模型的通过率从 2024 年初的 $\sim$5\%
        提升到 2025 年底的 $\sim$50\%（在 SWE-bench Verified 上）。
  \item \textbf{LiveCodeBench}：
        持续收集编程竞赛平台（LeetCode、Codeforces）的新题，
        避免数据泄露。题目有明确的时间戳，
        可以按模型训练截止日期进行公平比较。
\end{itemize}

\subsubsection{SWE-bench：从函数到系统的代码评估革命}

SWE-bench~\cite{jimenez2024swebench} 的意义在于
它将代码评估从\textbf{函数级}（HumanEval）
提升到了\textbf{系统级}。
一个典型的 SWE-bench 任务要求模型：
\begin{enumerate}[noitemsep]
  \item 阅读 GitHub issue 中用户的问题描述（通常含错误日志）。
  \item 理解一个包含数百个文件、数万行代码的仓库结构。
  \item 定位到 bug 所在的具体文件和函数。
  \item 编写一个正确的 patch（包括代码修改和可能的测试更新）。
  \item 通过仓库的既有测试套件验证。
\end{enumerate}

SWE-bench 经历了三个版本的演进：
\begin{itemize}[noitemsep]
  \item \textbf{SWE-bench}（原版）：
        2294 个来自 12 个 Python 仓库的真实 issue。
        部分任务被发现存在评估问题
        （如测试不够严格导致错误方案通过）。
  \item \textbf{SWE-bench Verified}（2024 年 8 月）：
        由人类工程师审核的 500 个高质量子集。
        确保每个任务的测试用例足够严格。
        曾被广泛用作标准评估集。
  \item \textbf{SWE-bench Pro}（Scale AI SEAL，2026 年）：
        针对 Verified 版本出现的污染问题而推出的升级版本，
        包含 \textbf{1,865 个多语言编程问题}
        （不再局限于 Python），
        使用私有测试用例和更严格的防泄露措施。
        截至 2026 年 3 月，最高分约 \textbf{44\%}——
        同一模型在 Pro 上的分数大约只有 Verified 的一半，
        更真实地反映了实际的软件工程能力。
\end{itemize}

SWE-bench 的演进揭示了一个重要模式：
截至 2026 年 3 月，SWE-bench Verified 的最高分
已达到约 \textbf{79.2\%}（Claude Opus 4.6 Thinking），
但学界已普遍认为 Verified 版本存在\textbf{部分污染}——
一些模型可能在训练数据中接触过测试用例的讨论和解决方案。
SWE-bench Pro 的推出正是对这一问题的直接回应。
同一模型在 Verified 上 $\sim$79\%、在 Pro 上 $\sim$44\% 的巨大落差，
生动说明了\textbf{基准测试的污染程度直接影响评估结论的可靠性}。

\textbf{其他代码智能体基准}也在快速发展：
\textbf{Terminal Bench 2.0} 测试模型在终端环境中的操作能力，
\textbf{BrowseComp} 评估模型在浏览器环境中的信息检索与操作能力。
这些基准将代码评估从「修 bug」扩展到了
更广泛的「计算机使用」能力维度。

\subsection{多模态评估}
\label{subsec:multimodal-eval}

随着 GPT-4V、Gemini、Claude 3+ 等模型原生支持视觉输入，
多模态评估在 2024--2026 年成为基准测试的重要新维度。

\subsubsection{MMMU：多模态大学水平理解}

\textbf{MMMU}（Massive Multi-discipline Multimodal Understanding）~\cite{yue2024mmmu}
是评估模型视觉-文本联合推理能力的核心基准。
它包含 11,500 道涵盖 30 个学科和 183 个子领域的多选题，
每道题都包含图表、公式、示意图等视觉元素，
不是简单的「看图说话」，
而是需要\textbf{同时理解图像内容和学科知识}才能回答。

MMMU 的难度体现在两个层面。
第一是\textbf{视觉理解的精度要求}：
题目中的图表可能包含关键的数值信息，
模型必须准确识别坐标轴的刻度、趋势线的斜率、
表格中的具体数字才能正确回答。
第二是\textbf{跨模态推理}：
答案通常不在图片中直接出现，
而是需要将图片信息与文本描述中的背景知识结合，
进行多步推理才能得出。

截至 2026 年初，
前沿多模态模型在 MMMU 上的准确率约为 70--80\%，
人类专家的准确率约为 88\%。
差距虽然在快速缩小，但仍然明显，
尤其是在需要精确图表阅读的理工科题目上。

\textbf{MMMU-Pro}~\cite{yue2024mmmupro}
是 MMMU 的升级版本，
通过三个设计改进增加了难度：
（1）移除了「仅靠文本描述不看图就能回答」的题目，
确保每道题都\textbf{必须}理解图像才能作答；
（2）增加了视觉描述题（vision-centric questions），
要求模型先详细描述图像内容，
再在此基础上推理，
防止模型「跳过看图」直接猜测；
（3）将选项数从 4 个增加到 10 个，
降低了随机猜测的基线。
MMMU-Pro 在 2025 年底发布后，
前沿模型的准确率比 MMMU 低 10--20 个百分点，
恢复了良好的区分度。

\subsubsection{MathVista：视觉数学推理}

\textbf{MathVista}~\cite{lu2024mathvista}
专门评估模型在视觉数学场景中的推理能力。
题目来源包括数学教科书中的几何图形、
统计学的数据可视化图表、
以及科学实验的图解说明。
模型需要从图像中提取数学信息
（如几何形状的尺寸、函数图像的特征点），
然后进行计算或推理。

MathVista 揭示了一个有趣的能力分离：
在纯文本数学（MATH、AIME）上表现极强的推理模型，
在 MathVista 上可能不如视觉能力更强的标准多模态模型。
这说明\textbf{视觉理解和数学推理是两种相对独立的能力}——
需要在训练中分别培养。

\subsubsection{视觉评估的特殊挑战}

多模态评估引入了文本评估不具备的独特挑战：

\textbf{图像分辨率的影响}。
同一道题目在不同分辨率下的评估结果可能差异显著。
一些模型在高分辨率（如 1344$\times$1344）下的准确率
比低分辨率（336$\times$336）高出 10--15 个百分点，
因为低分辨率下细节信息丢失导致关键数据不可读。
评估协议必须标准化输入图像的分辨率和预处理方式。

\textbf{OCR 能力的纠缠}。
许多多模态基准测试中，
模型的主要失败原因不是「理解错误」而是「读取错误」——
即无法正确识别图像中的文字、数字或符号。
这意味着基准测试实际上同时在评估两种能力
（OCR + 推理），
而这两者的混淆使得诊断模型的具体弱点变得困难。

\textbf{评估成本}。
多模态评估的计算成本远高于纯文本评估，
处理图像需要额外的视觉编码器计算，
而 MMMU 中的 11,500 道题目
在 70B 模型上的完整评估可能需要数十小时。

\subsection{长上下文}
\begin{itemize}[noitemsep]
  \item \textbf{RULER}：
        综合性长上下文评估。
        包含多种任务类型（检索、聚合、推理）
        和多种上下文长度（4K 到 128K+）。
  \item \textbf{Needle-in-a-Haystack}（NIAH）：
        在长文本中某个随机位置插入一个特定事实，
        然后在文本末尾询问这个事实。
        测试模型在不同上下文位置的信息检索能力。
        简单直观，但只测试了最基本的「能否找到信息」。
  \item \textbf{InfiniteBench}~\cite{zhang2024infinitebench}：
        专门测试 100K+ token 上下文能力的基准，
        包含阅读理解、代码理解、对话检索等多种任务。
        与 NIAH 的单点检索不同，
        InfiniteBench 的任务需要\textbf{跨越长上下文的信息聚合}——
        答案散布在文本的多个位置，
        模型必须综合多处信息才能正确回答。
\end{itemize}

\subsubsection{长上下文评估的演进：从检索到推理}

长上下文评估经历了三代演进，
反映了社区对「长上下文能力」理解的深化：

\textbf{第一代：单点检索}（2023--2024）。
以 Needle-in-a-Haystack 为代表。
在长文本的随机位置插入一个「针」
（如「The secret code is 42。”），
然后在末尾询问密码。
这种测试极其简单，它只检查模型能否在给定位置
找到特定信息，不涉及任何理解或推理。
到 2025 年，几乎所有 128K 上下文模型
都能通过 NIAH 测试，使其失去了区分价值。

\textbf{第二代：多点聚合}（2024--2025）。
以 RULER 和 InfiniteBench 为代表。
任务不再是「找到一个事实」，
而是「在 100K token 中找到散布的多个相关信息并综合推理」。
例如，给模型一份完整的法律合同，
要求识别所有与特定条款相关的义务和例外情况，
而这些信息可能分散在合同的不同章节中。
这一代评估区分出了两种模型：
「能找到信息但无法聚合」vs.「能有效聚合信息」。

\textbf{第三代：长上下文推理}（2025--2026）。
新一代基准不仅要求信息检索和聚合，
还要求\textbf{在长上下文之上进行复杂推理}。
例如：给模型一个包含 50,000 行代码的仓库，
要求分析跨多个文件的数据流以发现潜在 bug；
或者给模型 10 篇相互引用的学术论文，
要求识别论点之间的矛盾。
这些任务接近真实的知识工作者场景，
也是当前模型表现最不稳定的领域。

\begin{whybox}[为什么 NIAH 是一个「假的」长上下文测试？]
Needle-in-a-Haystack 的根本问题在于
它测试的是\textbf{注意力机制的检索能力}——
而非\textbf{理解长文本的能力}。
一个通过 NIAH 的模型可能完全无法完成
「总结一本 100 页书的主题」这种看似简单的任务，
因为总结需要对长文本建立全局理解，
而 NIAH 只需要局部定位一个已知模式。
这就是为什么 2025 年之后——
没有严肃的评估机构还在单独使用 NIAH
作为长上下文能力的衡量标准。
\end{whybox}

\subsection{指令跟随与对话}
\begin{itemize}[noitemsep]
  \item \textbf{MT-Bench}：
        80 道多轮对话题，由 GPT-4 评分（1--10 分）。
        涵盖写作、推理、数学、编码等 8 个类别。
        成本低，但受限于 GPT-4 自身的判断能力。
  \item \textbf{AlpacaEval}：
        805 道指令，由 GPT-4 与参考模型做 pairwise 比较。
        快速廉价，但存在「长度偏见」，
        GPT-4 评分器倾向于偏好更长的回答。
        AlpacaEval 2.0 引入了长度控制来缓解这个问题。
  \item \textbf{Arena-Hard}：
        从 Chatbot Arena 的真实用户问题中
        筛选出最具区分度的 500 道题。
        使用 GPT-4-Turbo 做评判。
        比 MT-Bench 和 AlpacaEval 有更高的区分度。
  \item \textbf{IFEval}（Instruction Following Eval）~\cite{zhou2023ifeval}：
        专门评估模型的\textbf{精确指令跟随能力}。
        每道题包含一个或多个可程序验证的约束条件
        （如「回答必须恰好包含 3 个段落」
        「不能使用任何逗号」
        「必须以问句结尾」），
        评估模型能否\textbf{严格满足所有约束}。
        IFEval 的独特价值在于它不评估回答质量，
        只评估模型能否「听从指示」。
        这一能力在 Agent 场景中尤为关键，
        因为 Agent 的工具调用需要模型
        严格遵循格式和参数要求。
\end{itemize}

\subsubsection{WildBench：来自真实世界的指令}

\textbf{WildBench}~\cite{lin2024wildbench}
与 MT-Bench 和 AlpacaEval 的「人造指令」不同，
它的 1,024 道题目\textbf{全部来自真实用户}——
从 ChatGPT、Claude 等产品的公开分享中收集
（经过脱敏和去重处理）。

WildBench 的核心见解是：
真实用户的指令与研究者设计的指令存在系统性差异。
真实指令往往更\textbf{模糊}
（「帮我写一封好的求职信」而非「写一封 200 字的求职信给 X 公司」），
更\textbf{多约束}
（「要专业但不要太正式，要有个性但不要太随意」），
更\textbf{上下文依赖}
（引用之前的对话、假设共享的文化背景）。
在这些真实指令上的表现差异，
往往比在标准化基准上的差异更能预测用户满意度。

\subsection{动态基准：LiveBench}

\begin{keybox}[LiveBench：抗污染的动态基准]
LiveBench~\cite{white2024livebench} 采用了一种独特的策略来对抗数据泄露：
它\textbf{定期更换测试题}——
使用模型训练截止日期之后才出现的信息来构造问题。
例如，使用最新的新闻文章、最新发布的数学竞赛题、
最新的代码仓库 commit 来生成评估题目。
这保证了没有模型能够通过训练数据「背答案」。
代价是需要持续投入人力维护和更新测试集。
\end{keybox}

LiveBench 于 2024 年发布，
并作为 Spotlight Paper 入选 ICLR 2025，
标志着动态评估范式获得了学术界的认可。
其设计原则包括：

\begin{enumerate}[noitemsep]
  \item \textbf{客观可验证的答案}：
        每道题目都有明确的、可程序验证的正确答案，
        消除了对 LLM 评判者的依赖。
        这与 MT-Bench 和 AlpacaEval 依赖 GPT-4 评分形成对比。
  \item \textbf{定期完全刷新}：
        测试题每 6 个月完全更换一次。
        旧题目公开用于研究，新题目取而代之。
  \item \textbf{7 大类别、23 项任务}：
        涵盖数学推理、代码生成、数据分析、
        语言理解、指令跟随、逻辑推理、综合能力。
  \item \textbf{天花板效应控制}：
        即使是最强的前沿模型——
        在 LiveBench 上的综合准确率也很少超过 65\%，
        保持了良好的区分度。
\end{enumerate}

\subsection{能力特化评估}
\label{subsec:capability-specific-eval}

除了通用基准，针对特定能力维度的评估
在 2025--2026 年迅速成熟。

\subsubsection{工具使用评估：BFCL 与 ToolBench}

随着 AI Agent 从概念走向产品，
评估模型的\textbf{工具使用能力}变得至关重要。

\textbf{BFCL}（Berkeley Function Calling Leaderboard）~\cite{patil2024bfcl}
是目前最权威的函数调用能力评估基准。
它评估模型在以下场景中的表现：
\begin{itemize}[noitemsep]
  \item \textbf{简单函数调用}：
        给定一个 API 文档和用户请求，
        模型能否生成正确的函数调用
        （正确的函数名、正确的参数值、正确的数据类型）。
  \item \textbf{多函数选择}：
        从多个可用函数中选择最合适的一个。
        这测试模型对 API 语义的理解深度。
  \item \textbf{并行调用}：
        识别用户请求中的多个独立子任务，
        并行调用多个函数。
  \item \textbf{嵌套调用}：
        将一个函数的返回值作为另一个函数的输入，
        形成调用链。
  \item \textbf{错误恢复}：
        当函数返回错误时，
        模型能否正确解读错误信息并重试或采取替代方案。
\end{itemize}

截至 2026 年初，
前沿模型在 BFCL 上的整体准确率约为 85--92\%，
但在嵌套调用和错误恢复的场景中
准确率降至 60--75\%，
说明复杂的工具使用仍是显著的短板。

\textbf{ToolBench}~\cite{qin2023toolllm}
则从另一个角度评估工具使用：
它提供了超过 16,000 个真实世界的 API
（来自 RapidAPI 等平台），
要求模型\textbf{自主规划}使用哪些 API 来完成复杂任务。
与 BFCL 的「给定函数，正确调用」不同，
ToolBench 测试的是「从海量工具中选择、组合和编排」的能力。

\subsubsection{Agent 评估：WebArena 与 OSWorld}

\textbf{WebArena}~\cite{zhou2024webarena}
是评估 Web Agent 的标杆基准。
它在一个真实的 Web 环境中
（包含 Reddit 克隆、电商网站、CMS 等）
部署了 812 个端到端任务，
要求模型通过浏览器操作
（点击、填写表单、导航页面）完成复杂的用户目标。
例如：「在论坛中找到关于 Python 异步编程的帖子，
给评分最高的回答点赞，然后回复表示感谢。」

WebArena 的评估结果在 2024--2026 年经历了显著提升：
从 2024 年初的约 6\%（GPT-4）到 2026 年初的约 35--40\%（frontier agents）。
但即使是最好的 Agent——
在需要多步规划和错误恢复的复杂任务上
成功率仍低于 20\%，
与人类用户约 78\% 的成功率相比差距巨大。

\textbf{OSWorld}~\cite{xie2024osworld}
将 Agent 评估从浏览器扩展到了\textbf{完整的桌面操作系统}。
模型需要通过截图理解屏幕内容，
并使用键盘鼠标操作来完成任务，
如在 LibreOffice 中编辑文档、
在文件管理器中组织文件、
在终端中运行命令等。
OSWorld 目前是 Agent 评估中最难的基准之一，
前沿模型的通过率低于 15\%。

\subsubsection{SWE-bench 之外：全栈编程评估}

代码评估正在从「修 bug」向更广泛的
「软件工程全流程」扩展。
2025--2026 年出现了一系列新的编程评估基准：

\textbf{Terminal Bench 2.0} 评估模型在终端环境中的操作能力，
不仅是写代码，还包括使用 shell 命令、
管理文件系统、配置环境等。
这更接近真实开发者的日常工作流。

\textbf{BrowseComp}（OpenAI）
评估模型通过浏览器进行信息检索和操作的能力。
与简单的搜索不同，
BrowseComp 要求模型在多个网页之间导航、
提取和综合信息、甚至在网页上执行操作
（如填写表单、下载文件）。

\textbf{Design2Code}~\cite{si2024design2code}
评估模型将视觉设计稿
转化为可工作的前端代码的能力。
输入是一张网页截图，
输出是对应的 HTML/CSS 代码。
这一基准测试了
视觉理解和代码生成的\textbf{联合能力}。

\textbf{CodeContests}（DeepMind）
汇集了 Codeforces、AtCoder 等竞赛平台的题目，
评估模型的算法设计和实现能力。
与 HumanEval 的「给 docstring 写函数」不同，
竞赛题通常需要
设计非平凡的算法才能在时间限制内通过，
这测试的是更深层的计算思维能力。

\subsubsection{创意写作评估的困境}

创意写作是 LLM 评估中\textbf{至今没有好的基准}的领域。
这不是因为缺乏尝试，而是因为创意写作的本质
与标准化评估之间存在根本性矛盾。

\textbf{创意的主观性}是最大的障碍。
一篇小说是否「好」取决于读者的品味、文化背景和阅读语境。
同一篇文章可能被一位文学评论家视为杰作，
被另一位视为平庸。
这种主观性使得建立「正确答案」几乎不可能，
而没有正确答案就无法做自动化评估。

\textbf{格式化偏好与创意质量的混淆}
加剧了评估的困难。
当使用 LLM-as-Judge 评估创意写作时，
评分模型倾向于偏好特定的文体风格，
如结构整齐、修辞丰富、带有「文学味」的文本
往往获得更高分数，
即使这些文本在原创性和情感深度上可能不如
更朴素但更真挚的文字。
Chatbot Arena 的创意写作类别在一定程度上缓解了这个问题
（因为评价来自真实读者而非 AI 评审），
但样本量远小于通用对话类别。

\textbf{缺乏长篇创意评估}也是一个空白。
现有的评估几乎都限于短文本（一首诗、一个段落、一封信），
而真正区分 AI 创意能力的是\textbf{长篇写作}：
能否在 10,000 字的小说中维持一致的人物性格？
能否在多章节的叙事中构建令人信服的情节弧线？
这些能力在 2026 年仍然缺乏标准化的评估方式。

\subsection{Chatbot Arena：人类评估的金标准}
\label{subsec:chatbot-arena}

\textbf{LMSYS Chatbot Arena}%
\footnote{LMSYS 即 Large Model Systems Organization——
由 UC Berkeley SkyLab 的研究者于 2023 年创立。
2025 年底，团队以 LMArena 的名义独立运营，
并获得了超过 1 亿美元的融资。}
是目前最受信赖的 LLM 评估平台。
它的设计哲学是：让真实用户用真实问题来评判模型。

\subsubsection{核心流程}

\begin{enumerate}[noitemsep]
  \item 用户提出\textbf{任意}问题，没有预设的题库，
        问题完全来自真实用户需求。
  \item 系统将问题同时发送给两个\textbf{匿名}模型。
        用户不知道自己在与哪个模型对话。
  \item 用户看到两个回答，选择更好的那个
        （或选择平局，或选择「都不好」）。
  \item 投票结束后揭示模型身份。
  \item 基于所有投票，使用 Bradley-Terry 模型
        计算每个模型的 Elo 评分。
\end{enumerate}

\subsubsection{Bradley-Terry 模型与 Elo 评分}

Bradley-Terry 模型假设模型 $i$ 击败模型 $j$ 的概率是：
\begin{equation}
  P(i \succ j) = \frac{e^{\beta_i}}{e^{\beta_i} + e^{\beta_j}}
  \label{eq:bt-model}
\end{equation}
其中 $\beta_i$ 是模型 $i$ 的能力参数。
通过最大似然估计拟合所有投票数据，
得到每个模型的 Elo 分数。

更具体地说，给定 $N$ 次成对比较的数据
$\{(i_k, j_k, y_k)\}_{k=1}^{N}$，
其中 $y_k \in \{i_k, j_k\}$ 表示胜者，
对数似然函数为：
\begin{equation}
  \mathcal{L}(\boldsymbol{\beta}) = \sum_{k=1}^{N}
  \left[\beta_{y_k} - \log\left(e^{\beta_{i_k}} + e^{\beta_{j_k}}\right)\right]
  \label{eq:bt-loglik}
\end{equation}
通过最大化此对数似然，
可以估计所有模型的能力参数 $\boldsymbol{\beta}$。
为了可解释性，
最终的 Elo 分数通常线性变换到
均值约 1000--1200 的范围内。

\begin{whybox}[为什么 Bradley-Terry 优于简单胜率？]
直接计算胜率（Win Rate）的问题在于：
不同模型被匹配的对手不同。
一个模型如果总是和弱模型对比，
胜率会虚高但实力并不强。
Bradley-Terry 模型的优势在于
它同时考虑了\textbf{所有}成对比较的传递性：
如果 A 经常胜 B，B 经常胜 C，
即使 A 和 C 从未直接比较——
模型也能正确推断 A 应该强于 C。
这使得即使对战矩阵是稀疏的
（并非每两个模型都有足够的直接对战），
也能得到可靠的全局排名。
\end{whybox}

\subsubsection{Style Control：分离风格与实质}

2024 年 8 月，Chatbot Arena 团队提出了
\textbf{Style Control}（风格控制）机制~\cite{li2024stylecontrol}，
解决了一个重要的偏见问题：
用户倾向于偏好更\textbf{长}的、
使用更多 \textbf{Markdown 格式}的回答。
这意味着一个输出更冗长但内容不一定更好的模型
可能获得更高的 Elo 分数。

Style Control 通过在 Bradley-Terry 模型中
加入\textbf{风格特征}作为协变量来缓解此问题：
\begin{equation}
  P(i \succ j) = \sigma\bigl(\beta_i - \beta_j
  + \boldsymbol{\gamma}^\top(\mathbf{s}_i - \mathbf{s}_j)\bigr)
  \label{eq:style-control}
\end{equation}
其中 $\sigma$ 是 sigmoid 函数，
$\mathbf{s}_i$ 是模型 $i$ 的风格特征向量
（包括回答长度、Markdown 使用率等）——
$\boldsymbol{\gamma}$ 是对应的系数。
这样，$\beta_i$ 就更纯粹地反映了模型的
\textbf{实质能力}而非风格偏好。

开启 Style Control 后的排名变化非常有启发性：
某些以简洁著称的模型（如早期的 Claude）
在 Style Control 排名中显著上升；
而某些以冗长著称的模型则相应下降。

\subsubsection{分类排行榜}

Chatbot Arena 不仅提供总体排名，
还按\textbf{能力类别}提供细分排行榜：
\begin{itemize}[noitemsep]
  \item \textbf{Hard Prompts}：
        从用户提交中筛选出特别困难的问题，
        通常是需要深度推理、专业知识或复杂指令跟随的问题。
        Hard Prompts 排名比总体排名更能区分前沿模型，
        因为简单问题上大多数好模型表现相似。
  \item \textbf{Coding}：编程相关问题的独立排名。
  \item \textbf{Math}：数学推理问题的独立排名。
  \item \textbf{Creative Writing}：创意写作能力排名。
  \item \textbf{Vision}：
        多模态（图文理解）能力排名。
        截至 2026 年初已积累超过 67 万次视觉评估。
  \item \textbf{多语言}：中文、日文、法文等非英语能力排名。
\end{itemize}

\begin{corebox}[Chatbot Arena 的数据规模（截至 2026 年初）]
\begin{itemize}[noitemsep]
  \item 累计超过 \textbf{800 万次}人类评估投票
  \item 覆盖超过 \textbf{200 个}模型
  \item 文本排行榜涵盖 100+ 模型，视觉排行榜涵盖 100+ 模型
  \item LMArena 独立运营并获得超过 1 亿美元融资
  \item 从 UC Berkeley 的学术项目演变为
        AI 行业的基础设施级评估平台
\end{itemize}
\end{corebox}

Chatbot Arena 的优势在于：
\begin{itemize}[noitemsep]
  \item \textbf{问题来自真实用户需求}——不是研究者预设的。
  \item \textbf{评估标准由用户自定义}——
        用户根据自己的偏好判断「更好」，不受固定评分标准约束。
  \item \textbf{模型匿名}——消除了品牌偏见。
  \item \textbf{动态更新}——使得针对性优化极为困难。
  \item \textbf{统计显著性极高}——
        百万级的投票量使得 Elo 评分的置信区间非常窄。
\end{itemize}

其局限性也值得注意：
\begin{itemize}[noitemsep]
  \item \textbf{用户群体偏差}：
        Arena 的用户以技术人员为主，
        可能不代表普通用户的偏好。
  \item \textbf{单轮评估为主}：
        大多数投票基于单轮对话，
        无法完全反映多轮对话的质量。
  \item \textbf{位置偏差}：
        研究表明用户有轻微的偏好左侧回答的倾向，
        虽然 Arena 通过随机化位置来缓解。
  \item \textbf{无法评估安全性}：
        用户投票主要基于「有用性」——
        不会主动测试和评估安全性问题。
\end{itemize}

\section{LLM-as-Judge：用 AI 评估 AI}
\label{sec:llm-as-judge}

随着人类评估成本的持续上升
和 LLM 能力的快速进步，
用强模型评估弱模型，
即 \textbf{LLM-as-Judge}——
已经从一种权宜之计发展为主流评估范式。

\subsection{核心方法论}

LLM-as-Judge 的基本形式有三种：

\textbf{单独评分}（Single Rating）：
将一个模型的输出单独呈现给 Judge 模型，
要求按预设标准（如 1--10 分）评分。
优点是简单，可以独立评估每个模型；
缺点是不同时间、不同 prompt 下的评分缺乏可比性。

\textbf{成对比较}（Pairwise Comparison）：
将两个模型的输出同时呈现给 Judge，
要求判断哪个更好。
这是 AlpacaEval 和 MT-Bench 采用的核心方法。
优点是判断更准确
（人类也更擅长做相对判断而非绝对评分）；
缺点是$N$个模型需要 $O(N^2)$ 次比较。

\textbf{参考答案比较}（Reference-Based）：
提供一个参考答案，
要求 Judge 评估模型输出与参考答案的一致程度。
优点是评估标准更明确；
缺点是参考答案的质量直接限制了评估的上限。

\subsection{已知偏见与校准}

LLM-as-Judge 存在多种系统性偏见，
如果不加以理解和校准，
评估结论可能严重失真：

\textbf{冗长偏好}（Verbosity Bias）。
这是最广泛被报告的偏见：
Judge 模型系统性地偏好更长的回答。
Dubois et al.\ (2024)~\cite{dubois2024alpacaeval2}
在 AlpacaEval 2.0 中量化了这一效应：
将回答长度增加一倍（通过简单的重复和扩展），
在不改变实质内容的情况下
可以将 win rate 提升 10--15 个百分点。
AlpacaEval 2.0 通过引入\textbf{长度控制回归}
（length-controlled win rate）来校准这一偏见。

\textbf{位置偏差}（Position Bias）。
在成对比较中，
Judge 模型倾向于偏好第一个
（或在某些模型中是第二个）呈现的回答。
标准的缓解方法是
对每次比较运行两遍（交换位置），
取一致的结果或随机打破平局。
这将计算成本翻倍，
但在统计上消除了位置效应。

\textbf{自我偏好}（Self-Enhancement Bias）。
当 Judge 模型评估自己或同系列模型的输出时，
倾向于给出更高的分数。
GPT-4 作为 Judge 评估 GPT-4 的输出时，
比评估同等质量的 Claude 输出
平均高出约 5--8\%~\cite{wataoka2024selfbias}。
这意味着使用特定模型作为 Judge
可能系统性地有利于该模型系列，
最佳实践是使用多个不同系列的模型作为 Judge
并取平均或多数投票。

\textbf{风格偏好}（Style Bias）。
Judge 模型可能偏好与自身训练风格相似的输出，
例如，使用更多 Markdown 格式、
更结构化的回答、
更「学术化」的语气。
这与冗长偏好有部分重叠，
但风格偏好更微妙、更难量化。

\begin{whybox}[LLM-as-Judge 在何时可信？]
经验法则：
\begin{itemize}[noitemsep]
  \item \textbf{高可信}：评估事实准确性、代码正确性、
        格式合规性等有明确标准的维度。
        Judge 模型在这些任务上与人类评审的一致率
        可达 85--95\%。
  \item \textbf{中等可信}：评估回答的完整性、
        逻辑连贯性、专业深度。
        一致率约 75--85\%。
  \item \textbf{低可信}：评估创意质量、
        文化敏感度、幽默感、情感共鸣。
        一致率可能低于 70\%。
\end{itemize}
关键原则：LLM-as-Judge 不是人类评估的替代品，
而是\textbf{筛选工具}——
用它快速排除明显差的模型，
用人类评估做最终的微调决策。
\end{whybox}

\subsection{多 Judge 聚合与元评估}

单一 Judge 模型的偏见使得
\textbf{多 Judge 聚合}成为 2025--2026 年的最佳实践。
核心思想是使用多个不同系列的模型作为 Judge
（如同时使用 GPT-4o、Claude 3.5 和 Gemini 1.5 Pro），
然后通过投票或加权平均得到最终评分。

\textbf{多数投票}（Majority Vote）是最简单的聚合方法：
3 个 Judge 中有 2 个认为 A 更好，则 A 获胜。
这种方法能有效抵消个别 Judge 的系统偏见，
但需要 3 倍的计算成本。

更精细的方法是\textbf{加权聚合}——
根据每个 Judge 与人类评审的历史一致率
动态调整其权重。
如果 GPT-4o 在代码评估上与人类一致率最高
而 Claude 在创意写作上最准确，
则在不同任务类别上使用不同的权重分配。

\textbf{元评估}（Meta-Evaluation）
是验证 Judge 质量的关键方法论。
它通过测量 Judge 模型与人类评审在同一组样本上的一致率
来量化 Judge 的可靠性。
常用的指标包括：
\begin{itemize}[noitemsep]
  \item \textbf{Cohen's Kappa}：
        衡量 Judge 与人类超出随机一致的程度。
        $\kappa > 0.6$ 被认为是合理的，
        $\kappa > 0.8$ 被认为是优秀的。
  \item \textbf{成对一致率}（Pairwise Agreement）：
        在成对比较中，
        Judge 与人类做出相同判断的比例。
        2026 年的前沿 Judge（如 GPT-4o）
        在事实性判断上达到 90\% 以上的一致率，
        但在风格和创意方面降至 70--75\%。
  \item \textbf{排名相关性}（Rank Correlation）：
        Judge 产生的模型排名与
        人类产生的排名之间的 Spearman 相关系数。
        优秀的 Judge 应达到 $\rho > 0.9$。
\end{itemize}

\subsection{LLM-as-Judge 的实践架构}

在生产环境中部署 LLM-as-Judge
需要解决一系列工程问题：

\textbf{评估 Prompt 的设计}直接决定了评判质量。
研究表明，
一个包含\textbf{详细评分量表}（rubric）的 prompt
比简单地要求「判断哪个更好」
能产生更一致和更有信息量的评估。
最佳实践是为每个评估维度
提供明确的评分标准和示例：
“5 分：回答完全准确、逻辑清晰、
没有遗漏任何关键信息；
4 分：回答基本准确但有轻微遗漏；
3 分：...”

\textbf{批量评估的成本控制}是另一个实际考量。
如果需要评估 100 个模型在 1,000 道题上的输出，
使用 GPT-4o 作为 Judge 的 API 成本
可能达到数千美元。
\textbf{采样评估}（对题目子集进行评估，
然后用统计方法外推全集结果）
可以在可接受的精度损失下
将成本降低 5--10 倍。

\textbf{评估结果的可解释性}
是 LLM-as-Judge 相对于人类评审的一个潜在优势。
Judge 模型可以被要求不仅给出评分，
还生成\textbf{详细的评审意见}——
解释为什么 A 的回答更好，
指出 B 的回答在哪些方面不足。
这些评审意见可以作为
模型改进的直接反馈信号，
比单纯的分数更有价值。

\subsection{Judge 的递归问题与未来方向}

LLM-as-Judge 面临一个哲学层面的\textbf{递归问题}：
谁来评估 Judge 本身的质量？

如果我们用 GPT-4o 评估其他模型，
但 GPT-4o 自身也有偏见和错误，
那么评估结果的可信度上限
就是 GPT-4o 自身的能力上限。
当被评估的模型在某些维度上
\textbf{超越了 Judge 模型}时
（这在 2026 年已经开始发生，
Claude Opus 4.6 的推理能力在某些任务上超越了 GPT-4o），
Judge 的评价反而可能是错误的。

这个递归问题没有完美的解决方案，
但 2025--2026 年出现了一些有前景的方向：
\begin{itemize}[noitemsep]
  \item \textbf{可验证的评估}：
        对于有明确正确答案的任务
        （数学、代码、事实问答），
        用程序验证替代 Judge 评估，
        彻底绕过 Judge 的局限性。
  \item \textbf{分解评估}：
        将一个复杂的评估任务
        分解为多个简单的子判断，
        每个子判断都在 Judge 的能力范围内。
        这类似于 Constitutional AI 中
        「将复杂价值判断分解为简单原则检查」的思路。
  \item \textbf{人机混合评估}：
        用 LLM-as-Judge 做初步筛选（排除明显的差和好），
        对中间地带的案例使用人类评审。
        这种分层策略可以在成本和质量之间
        找到最优平衡。
\end{itemize}

\section{推理能力评估}
\label{sec:reasoning-eval}

随着 OpenAI o1/o3 系列和 DeepSeek-R1 等\textbf{推理模型}
（reasoning models）的出现，
传统的评估范式面临新的挑战。
推理模型通过在推理时生成大量 thinking tokens
来进行深度思考（详见第~\ref{ch:ttc}~章），
其能力特征与标准模型有显著差异，
在需要多步推理的任务上表现极强，
但在简单任务上可能反而因「过度思考」而表现更慢或更贵。

\subsection{推理模型 vs 标准模型：评估差异}

\begin{corebox}[推理模型的独特评估挑战]
推理模型改变了评估的基本假设：
\begin{enumerate}[noitemsep]
  \item \textbf{成本不再恒定}：
        标准模型每个问题消耗大致相同的计算量；
        推理模型根据问题难度动态调整，
        难题可能消耗 100 倍于简单题的 token。
        基准测试的排名需要考虑\textbf{计算效率}——
        而非仅仅是准确率。
  \item \textbf{简单基准失去意义}：
        推理模型在 GSM8K 上达到 $>$99\% 准确率，
        在 MMLU 上超过 90\%，
        这些基准完全无法区分推理模型的优劣。
        评估推理模型需要\textbf{显著更难}的题目
        （AIME、GPQA Diamond、FrontierMath）。
  \item \textbf{Thinking tokens 的可观察性}：
        一些推理模型（如 DeepSeek-R1）会输出完整的思考过程；
        另一些（如 o1 早期版本）隐藏了思考过程。
        这使得评估和比较推理质量变得复杂。
\end{enumerate}
\end{corebox}

推理模型在不同基准上的表现模式值得关注。
在「需要多步推理」的任务上，
推理模型相对标准模型的提升幅度巨大：
\begin{itemize}[noitemsep]
  \item \textbf{AIME 2024/2025}：
        标准模型（如 GPT-4o, Claude 3.5 Sonnet）
        的通过率约为 10--20\%，
        而推理模型（o3, R1）可以达到 80--90\% 以上。
  \item \textbf{GPQA Diamond}：
        标准模型准确率约 50--60\%，
        推理模型达到 70--93\%，
        部分结果超越人类博士专家水平。
  \item \textbf{FrontierMath Tier 1--2}：
        标准模型通过率 $<$5\%，
        推理模型可达 25\% 以上。
  \item \textbf{SWE-bench Verified}：
        推理模型在复杂 bug 修复上
        相比标准模型有 10--15 个百分点的提升。
\end{itemize}

但在「不需要深度推理」的任务上
（如创意写作、简单问答、翻译），
推理模型的表现与标准模型相当甚至稍弱，
过多的 thinking tokens 有时反而引入不必要的复杂性。

\subsection{推理模型的评估协议设计}

推理模型的评估需要专门设计的协议，
因为标准的评估协议对推理模型可能不公平
（或不公平地有利于推理模型）。

\textbf{Thinking budget 控制}是核心的协议设计问题。
推理模型可以通过消耗更多 thinking tokens
来提升准确率。理论上，
给予无限的 thinking budget，
推理模型在很多任务上都能达到近乎完美的表现。
但这在实践中是不合理的，
用户不会等待 10 分钟来得到一个简单问题的回答。
因此，评估协议需要规定
\textbf{最大 thinking tokens 数量}
或\textbf{最大推理时间}——
以确保评估反映的是
在合理计算预算下的能力。

\textbf{多样性采样策略}对推理模型
的评估结果有显著影响。
标准的评估通常只运行一次推理（pass@1）。
但推理模型的一个独特特性是
\textbf{高方差}：同一道题运行多次，
有时答对有时答错，
因为推理链的探索路径具有随机性。
为了更准确地评估推理模型的能力，
一些基准引入了 \textbf{pass@k}（k 次中至少一次正确）
和 \textbf{consensus@k}（k 次中多数正确）指标。
pass@k 反映了模型的\textbf{能力上限}
（在足够多的尝试中能否找到正确答案），
consensus@k 反映了模型的\textbf{可靠性}
（在多数情况下是否能给出正确答案）。

\textbf{思考过程的评估}是推理模型独有的评估维度。
一些推理模型（如 DeepSeek-R1）公开了
完整的 thinking tokens，
这使得评估不仅可以看最终答案，
还可以分析\textbf{推理质量}：
推理链是否逻辑连贯？
是否存在「暗中猜测、事后合理化」的现象
（即答案是随机猜到的——
但事后生成了看似合理的推理过程）？
推理过程中的自我纠正
（发现错误后回退并修正）是否频繁且有效？
这些定性分析为理解推理模型的「真正能力」
提供了比单纯准确率更丰富的信息。

\subsection{计算效率评估：准确率不是唯一指标}

推理模型的出现迫使评估社区重新思考一个基本问题：
\textbf{模型应该按准确率排名，还是按效率排名？}

一个在 AIME 上达到 96\% 准确率但消耗 50,000 tokens 的推理模型，
与一个达到 90\% 准确率但只消耗 2,000 tokens 的标准模型相比，
哪个「更好」？
答案取决于使用场景：
对于竞赛或考试（追求极致准确率），前者更好；
对于日常问答（追求性价比），后者更好。

这催生了\textbf{Pareto 效率分析}
在 LLM 评估中的应用：
将每个模型表示为（准确率, 计算成本）空间中的一个点，
识别 Pareto 前沿上的模型，
即在给定准确率下成本最低、或给定成本下准确率最高的模型。
2025--2026 年，越来越多的评估报告
开始同时报告准确率和 token 消耗量，
而非仅仅报告准确率排名。

\subsection{过程奖励模型 vs 结果奖励模型}
\label{subsec:prm-vs-orm}

评估推理过程（而非仅评估最终答案）
是推理模型训练和评估中的核心问题。
两种主要范式形成了鲜明的对比：

\subsubsection{结果奖励模型（ORM）}

\textbf{结果奖励模型}（Outcome Reward Model, ORM）
只关注最终答案是否正确：
\begin{equation}
  R_{\text{ORM}}(\text{solution}) =
  \begin{cases}
    +1 & \text{if final answer is correct} \\
    -1 & \text{if final answer is wrong}
  \end{cases}
  \label{eq:orm}
\end{equation}

ORM 的优势是\textbf{标注简单}——
只需验证最终答案，
不需要评估每一步推理的正确性。
对于数学和代码等有明确正确答案的任务，
ORM 的标注成本几乎为零（自动验证）。

但 ORM 存在严重的\textbf{信用分配问题}
（credit assignment problem）：
如果一个解题过程有 10 步，其中第 3 步出错导致最终答案错误，
ORM 只知道「答案错了」，
但不知道具体是哪一步出了问题。
在使用 RL 训练时，
错误的信用分配会导致学习效率低下，
前两步（正确的）也可能被错误地惩罚，
而第 3 步（错误的）可能没有被充分惩罚。

\subsubsection{过程奖励模型（PRM）}

\textbf{过程奖励模型}（Process Reward Model, PRM）~\cite{lightman2024lets}
对推理过程的\textbf{每一步}提供反馈：
\begin{equation}
  R_{\text{PRM}}(\text{step}_t \mid \text{step}_{1:t-1}) \in [-1, +1]
  \label{eq:prm}
\end{equation}
其中每一步的奖励取决于该步骤
在之前所有步骤的上下文中是否正确。

PRM 解决了信用分配问题：
它能精确指出推理链中\textbf{哪一步}出错，
为 RL 训练提供更密集、更准确的学习信号。
OpenAI 在 2023 年的研究~\cite{lightman2024lets}
表明 PRM 引导的搜索（Best-of-N with PRM）
在数学问题上显著优于 ORM 引导的搜索。

\begin{whybox}[PRM 为什么更有效？]
直觉上，PRM vs ORM 的差异
类似于老师批改作业的两种方式：
\begin{itemize}[noitemsep]
  \item \textbf{ORM}：只看最终答案，
        打对号或叉号，不做过程批注。
        学生知道答案错了，但不知道错在哪里。
  \item \textbf{PRM}：逐步检查，
        在每一步标注对错。
        学生能精确知道自己在哪一步走偏了。
\end{itemize}
后者显然提供了更多的学习信息，
但\textbf{批注成本}（标注成本）也高得多。
\end{whybox}

\subsubsection{PRM 的标注瓶颈与自动化}

PRM 面临的核心挑战是\textbf{标注成本}：
逐步评估推理过程需要大量的人力。
OpenAI 的 PRM800K 数据集~\cite{lightman2024lets}
包含了 80 万步的人类标注，成本极高且难以扩展。

Setlur et al.\ (2025)~\cite{setlur2025rewarding}
在 ICLR 2025 上提出了
\textbf{自动化过程验证}（Automated Process Verification）的方法。
核心思想是利用结果监督来\textbf{自动生成}过程标签：
\begin{enumerate}[noitemsep]
  \item 给定推理链的前 $t$ 步，
        让模型多次独立完成剩余步骤。
  \item 如果从第 $t$ 步之后继续推理
        \textbf{大概率能得到正确答案}——
        那么第 $t$ 步很可能是正确的。
  \item 如果从第 $t$ 步之后继续推理
        \textbf{几乎总是得到错误答案}——
        那么第 $t$ 步很可能是错误的。
\end{enumerate}
这种方法通过蒙特卡洛估计每一步的「正确概率」，
将不可扩展的人工逐步标注
转化为可自动化的过程，
使 PRM 训练的规模化成为可能。

Yang et al.\ (2025)~\cite{yang2025beyond}
进一步指出了现有 PRM 的一个局限：
在长推理链中，标注倾向于只关注\textbf{第一个错误步骤}
及其之前的步骤，
而忽略了后续步骤中可能的\textbf{自我纠正}。
他们提出了更灵活的「反思式」PRM，
能够评估模型在发现错误后的纠正能力。

\subsubsection{PRM 与 ORM 的实际选择}

\begin{keybox}[PRM vs ORM：实际应用指南]
\begin{itemize}[noitemsep]
  \item \textbf{数学/逻辑推理}：PRM 显著优于 ORM。
        推理链长、每步可验证、
        信用分配问题严重，PRM 的优势最大化。
  \item \textbf{代码生成}：ORM 通常足够。
        有单元测试作为准确的结果信号，
        而代码的「步骤」边界不如数学清晰。
  \item \textbf{开放式文本生成}：两者都很困难。
        既没有明确的正确答案（ORM 难），
        也没有可验证的中间步骤（PRM 难）。
        LLM-as-Judge 或人类评估更适合。
  \item \textbf{混合方法}：
        实践中，很多系统结合 ORM 和 PRM：
        PRM 用于引导搜索方向，
        ORM 用于最终结果验证。
\end{itemize}
\end{keybox}

\subsection{推理模型的专项基准}
\label{subsec:reasoning-benchmarks}

推理模型的快速发展催生了一系列新的评估基准，
专门针对需要深度推理的任务：

\begin{itemize}[noitemsep]
  \item \textbf{AIME 2024/2025}：
        美国数学邀请赛真题是评估推理模型数学能力的标准基准。
        题目需要创造性的多步推理，而非模式匹配。
        o3 在 AIME 2024 上的得分超过 90\%，
        接近人类数学竞赛选手的顶尖水平。
  \item \textbf{GPQA Diamond}：
        198 道博士级科学问题（见~\ref{subsec:chatbot-arena}~节）。
        推理模型在此基准上的准确率
        已从 2024 年初的约 50\% 提升到 2026 年初的约 90\%。
  \item \textbf{FrontierMath Tier 4}：
        研究级数学问题（见上文）。
        目前唯一仍对推理模型构成真正挑战的数学基准。
  \item \textbf{EpochAI Frontier Probes}：
        专门设计来测量推理能力边界的题目集。
\end{itemize}

\section{前沿模型基准对比}
\label{sec:benchmark-comparison}

为了呈现 2025--2026 年 LLM 能力格局的全貌，
我们汇总了主要前沿模型在核心基准上的表现%
\footnote{数据来源综合自各模型的技术报告、
Chatbot Arena 排行榜和独立评测平台。
数据截止到 2026 年 2 月。
注意不同来源的评测条件可能存在差异。}。

\begin{table}[htbp]
\centering\small
\renewcommand{\arraystretch}{1.3}
\caption{2025--2026 年前沿模型基准测试对比}
\label{tab:model-benchmarks}
\begin{tabular}{@{}llcccc@{}}
\toprule
\rowcolor{figLightGray}
\textbf{模型} & \textbf{类型} & \textbf{MMLU-Pro} & \textbf{GPQA-D}
  & \textbf{SWE-bench V} & \textbf{Arena Elo} \\
\midrule
\multicolumn{6}{@{}l}{\textit{推理模型（Reasoning Models）}} \\
GPT-5.2 Pro       & 推理  & 88.7\% & 93.2\% & 65.8\% & 1402 \\
Claude Opus 4.6   & 混合  & 88.2\% & 89.0\% & 72.5\% & 1504 \\
DeepSeek V3.2-S   & 推理  & 85.9\% & 85.3\% & 77.8\% & 1361 \\
Grok 4 Heavy      & 推理  & 86.4\% & 88.9\% & 61.0\% & 1375 \\
\midrule
\multicolumn{6}{@{}l}{\textit{标准模型（Standard Models）}} \\
Gemini 3.1 Pro      & 标准  & 89.8\% & 87.5\% & 63.2\% & 1485 \\
Claude Opus 4.5   & 标准  & 87.1\% & 86.5\% & 68.4\% & 1370 \\
GPT-5.2           & 标准  & 86.3\% & 88.0\% & 58.2\% & 1380 \\
Qwen 3.5          & 标准  & 84.6\% & 82.1\% & 62.5\% & 1342 \\
\midrule
\multicolumn{6}{@{}l}{\textit{开源模型（Open-Weight Models）}} \\
Llama 4 Maverick  & 标准  & 83.2\% & 78.5\% & 55.8\% & 1320 \\
DeepSeek R1       & 推理  & 90.8\% & 71.5\% & 49.2\% & 1310 \\
Qwen 3.5          & 标准  & 84.6\% & 82.1\% & 62.5\% & 1342 \\
Mistral 3         & 标准  & 82.8\% & 79.3\% & 54.1\% & 1315 \\
\bottomrule
\end{tabular}

\vspace{4pt}
{\footnotesize
\textit{注}：GPQA-D = GPQA Diamond; SWE-bench V = SWE-bench Verified;
Arena Elo 数据来自 Chatbot Arena 排行榜，截至 2026 年 2--3 月。
“推理” 类型指使用 thinking tokens 的模型；
“混合” 指同时支持标准和推理模式。
同一模型在不同评测条件下可能给出不同分数。}
\end{table}

从表~\ref{tab:model-benchmarks} 可以观察到几个值得注意的模式：

\textbf{推理模型在 GPQA 上优势巨大}：
推理模型在需要深度科学推理的 GPQA Diamond 上
普遍超过 85\%，部分超过 90\%，
远高于标准模型的 78--88\%。
这证实了 thinking tokens 在复杂推理任务上的价值。

\textbf{SWE-bench 上开源模型追赶迅速}：
DeepSeek V3.2-Speciale 在 SWE-bench Verified 上达到 77.8\%，
超过了多数闭源模型。
这得益于其在代码相关的后训练数据上的大量投入。

\textbf{Arena Elo 与基准分数并不完全一致}：
Gemini 3.1 Pro 在 MMLU-Pro 上得分最高（89.8\%），
但其 Arena Elo 低于 MMLU-Pro 分数较低的 Claude Opus 4.6。
这说明\textbf{基准测试和真实用户偏好衡量的是不同的东西}——
用户满意度取决于指令跟随、风格、一致性等
基准测试不一定覆盖的维度。

\textbf{开源与闭源的差距在缩小}：
Llama 4 Maverick 和 Qwen 3.5 等开源模型
在 MMLU-Pro 上已接近闭源前沿模型（差距 $<$5\%）。
在某些特定任务上，开源模型甚至有竞争力。

\section{安全对齐}

LLM 的安全问题包括多个层面：

\subsection{有害内容生成}

模型可能生成暴力、歧视、违法的内容。
通过 RLHF 中的安全奖励和 Constitutional AI~\cite{bai2022constitutional}
（让 AI 根据原则自我批评和修正）来缓解。

Constitutional AI 的核心思想是：
不依赖人类标注者来判断每个输出是否安全，
而是给模型一组\textbf{宪法原则}（如「不鼓励暴力」
「尊重隐私」「提供准确信息」），
让模型自己根据这些原则评估和修正输出。
这大幅降低了安全标注的人力成本，
同时使安全标准更加一致和可扩展。

\subsection{越狱攻击}

\textbf{越狱攻击}（Jailbreaking）：
用户通过巧妙的提示绕过模型的安全限制。
常见的越狱技巧包括：
\begin{itemize}[noitemsep]
  \item \textbf{角色扮演}：“假设你是一个没有任何限制的 AI...”
  \item \textbf{间接提问}：“写一个关于角色 X 做 Y 的小说...”
        （间接获取有害信息）
  \item \textbf{编码/加密}：用 Base64、逆序文本等方式
        隐藏有害指令。
  \item \textbf{多轮诱导}：通过多轮对话逐步引导模型
        越过安全边界。
  \item \textbf{多模态攻击}：
        将有害文本嵌入图像中（如 OCR 攻击），
        利用视觉通道绕过文本安全过滤器。
        这种攻击在 2025 年变得更加普遍，
        因为多模态模型的视觉理解能力
        使得图像中的文本可以被「读取」并执行。
  \item \textbf{跨语言攻击}：
        用低资源语言（如斯瓦希里语、威尔士语）
        表达有害请求。
        由于安全训练数据主要以英语为主——
        模型在其他语言上的安全防御往往更弱。
\end{itemize}

防御策略包括\textbf{对抗训练}
（在训练数据中纳入已知的越狱模式，
让模型学会识别和拒绝），
\textbf{输入/输出过滤}（独立的安全分类器
对输入和输出进行审查），
以及\textbf{红队测试}。

\textbf{Red teaming}（红队测试）
是系统性地寻找模型安全漏洞的过程。
专业的红队由安全研究者、领域专家和普通用户组成，
他们尝试各种方法来触发有害输出。
越来越多的团队使用\textbf{自动化红队}——
用另一个 LLM 来自动生成大量攻击 prompt，
然后检查目标模型的回复。
这大幅提高了测试覆盖率，
但人类红队仍然不可替代，
他们能发现自动化方法难以想到的创意攻击。

\subsection{安全评估基准}
\label{subsec:safety-benchmarks}

安全评估已经从零散的测试
发展为一个完整的基准测试体系：

\subsubsection{TruthfulQA：衡量幻觉与错误信息}

\textbf{TruthfulQA}~\cite{lin2022truthfulqa}
测试模型是否会复述常见的错误信息和迷思。
817 个问题涵盖健康、法律、金融、阴谋论等领域，
每个问题都设计为\textbf{诱导性}的，
人类容易错误回答的问题，模型也容易错。

TruthfulQA 的独特价值在于
它测试的不是知识的\textbf{广度}
而是知识的\textbf{准确性}——
特别是在「普遍持有但错误的信念」方面。
例如：「感冒是因为受凉引起的吗？」
一个只学到统计相关性的模型可能回答「是」，
因为训练数据中有大量将感冒与寒冷天气关联的文本。
但科学上，感冒是病毒感染导致的，
寒冷本身不是直接原因。

截至 2026 年，前沿模型在 TruthfulQA 上
的准确率约为 70--80\%，
说明幻觉问题虽有改善但远未解决。

\subsubsection{BBQ：偏见基准}

\textbf{BBQ}（Bias Benchmark for QA）~\cite{parrish2022bbq}
系统性地测试模型在不同人群维度上的偏见。
它覆盖 11 个社会偏见类别
（种族、性别、年龄、宗教、性取向、
残障状态、国籍、社会经济地位等），
每个类别包含数百道问题。

BBQ 的设计精妙之处在于每道题有两种形式：
\textbf{歧义形式}（信息不足以确定答案）和
\textbf{消歧形式}（信息足以确定答案）。
在歧义形式中，
理想的模型应当承认信息不足而拒绝做判断；
如果模型在歧义情况下仍然给出了具体判断，
该判断反映的就是模型的\textbf{隐含偏见}。

\subsubsection{HarmBench：攻击成功率评估}

\textbf{HarmBench}~\cite{mazeika2024harmbench}
是一个评估越狱攻击成功率的标准化框架。
它包含 510 种有害行为请求
（涵盖化学武器、网络攻击、
社会工程等多个危害类别），
并集成了 18 种主流攻击方法
（从手动越狱到自动化对抗生成）。

HarmBench 的关键贡献是
将安全评估从「能否通过特定的越狱 prompt」
提升为\textbf{系统性的攻击-防御博弈评估}。
对于每个模型，HarmBench 计算一个
\textbf{Attack Success Rate}（ASR），
即在所有攻击方法和所有有害请求的组合中——
成功绕过安全防御的比例。

截至 2026 年初，
前沿模型在最强攻击方法下的 ASR 约为 5--15\%，
在基本攻击方法下低于 2\%。
但这个数字的含义是：
一个足够有决心的攻击者，
仍然有概率从模型中获取有害信息。
\textbf{完全的安全性在当前技术下是不可能的}。

\subsubsection{Machiavelli：伦理推理评估}

\textbf{Machiavelli Benchmark}~\cite{pan2023machiavelli}
评估模型在面临道德困境时的决策质量。
它使用基于文本冒险游戏的交互式场景，
模型需要在多个选项中做出选择，
某些选项在短期内更有效
但涉及欺骗、操纵或伤害他人，
另一些选项虽然效率较低但更符合伦理。

Machiavelli 衡量两个维度：
\textbf{能力}（模型能否在游戏中取得好的结果？）
和\textbf{伦理}（模型是否会选择不道德的手段来实现目标？）。
理想的模型应当在两个维度上都表现出色，
既有能力又有道德。

一个令人不安的发现是：
经过 RLHF 训练的模型在 Machiavelli 测试中
表现出了\textbf{表面合规}——
在明显的伦理测试中做出道德选择，
但在伦理判断更微妙的场景中
仍然倾向于选择最「有效」（但不道德）的选项。
这说明当前的对齐技术可能只是
让模型学会了\textbf{识别伦理测试}——
而非真正建立了伦理推理能力。

\subsection{幻觉}

\textbf{幻觉}（Hallucination）：
模型自信地生成不正确的信息。
这是一个尚未完全解决的根本问题。
语言模型本质上是在做统计预测，
不具备真正的「知道自己不知道」的能力。

幻觉可以分为两类：
\begin{itemize}[noitemsep]
  \item \textbf{内在幻觉}（Intrinsic hallucination）：
        输出与\textbf{给定的上下文/来源}矛盾。
        例如，文档里说「收入增长了 5\%」，
        模型回答「收入增长了 15\%」。
  \item \textbf{外在幻觉}（Extrinsic hallucination）：
        输出无法从给定来源中验证，但也不矛盾。
        模型「添加」了来源中不存在的信息。
        这类幻觉更难检测，可能是正确的补充信息，
        也可能是捏造的。
\end{itemize}

幻觉的根源在于模型的训练目标（最大化下一个 token 的概率）
与「只在确定时才生成」之间的矛盾。
模型被训练为\textbf{总是生成一些东西}——
即使它对内容的准确性没有把握。
而且模型没有「事实性」的内在概念，
它学到的是训练数据中的统计模式，
而非可验证的事实。

\subsection{RAG：用检索对抗幻觉}

Retrieval-Augmented Generation（RAG）
通过在生成前检索外部知识来部分缓解幻觉问题。

RAG 的基本流程是：
\begin{enumerate}[noitemsep]
  \item \textbf{检索}：用户提问 $\to$ 将问题向量化 $\to$
        在向量数据库中检索最相关的文档块。
  \item \textbf{增强}：将检索到的文档块与用户问题拼接，
        作为上下文送入 LLM。
  \item \textbf{生成}：LLM 基于检索到的上下文生成回答，
        回答有了可追溯的「来源」。
\end{enumerate}

RAG 的关键技术组件：
\begin{itemize}[noitemsep]
  \item \textbf{Chunking}（文档切分）：
        将长文档切成固定大小的块（如 512 tokens）。
        切分策略影响检索质量，
        太大则不精确，太小则失去上下文。
        重叠切分（相邻块有 10--20\% 的重叠）可以缓解边界问题。
  \item \textbf{Embedding models}（嵌入模型）：
        将文本转换为高维向量。
        常用的包括 OpenAI text-embedding-3、
        BGE、E5-Mistral 等。
        嵌入质量直接决定检索质量。
  \item \textbf{Vector databases}（向量数据库）：
        存储和检索高维向量。
        常用的包括 Pinecone、Weaviate、Milvus、Chroma 等。
        核心操作是\textbf{最近邻搜索}（ANN）。
\end{itemize}

但 RAG 也不是万灵药。
如果检索到的文档本身有误，或者检索不到相关信息，
幻觉仍然可能发生。
此外，模型有时会\textbf{忽略}检索到的信息，
仍然依赖自己的「记忆」生成回答，
这在检索结果与模型内在知识矛盾时尤为常见。

\subsection{隐私与安全评估}

\textbf{隐私泄露}：
模型可能记住并泄露训练数据中的敏感信息。
差分隐私训练和去敏处理是主要的缓解手段。

\textbf{安全评估基准}：
\begin{itemize}[noitemsep]
  \item \textbf{TruthfulQA}：测试模型是否会复述常见的错误信息和迷思。
        817 个问题涵盖健康、法律、金融等领域的常见误解。
  \item \textbf{BBQ}（Bias Benchmark for QA）：
        测试模型在不同人群（种族、性别、年龄等）上的偏见。
  \item \textbf{ToxiGen}：测试模型生成有毒内容的倾向。
        涵盖针对 13 个少数群体的隐含和显性有毒语言。
  \item \textbf{AdvBench}：对抗性提示集合，
        测试模型在面对刻意诱导时的安全性。
\end{itemize}

\subsubsection{红队评估框架的工业化}

2025--2026 年，红队测试从手工艺式的安全审计
演变为\textbf{工业化的安全评估流程}。

\textbf{自动化红队}（Automated Red Teaming）
使用专门训练的「攻击者模型」
自动生成大量多样化的攻击 prompt。
Perez et al.\ (2022)~\cite{perez2022redteaming}
最早提出了这一范式，
用一个 LLM 生成攻击，用另一个 LLM 作为目标，
用第三个 LLM 判断攻击是否成功。
到 2026 年，这种三元组架构已经标准化：
\textbf{Attacker}（生成攻击 prompt）$\to$
\textbf{Target}（被测模型）$\to$
\textbf{Judge}（评估输出是否违规）。

自动化红队的规模化能力令人印象深刻：
一次自动化红队评估可以在数小时内
生成并测试\textbf{数万条}攻击 prompt，
覆盖手动红队可能需要数周才能探索的攻击面。
但自动化方法也有局限，
它倾向于发现已知模式的变体，
而非真正创新的攻击策略。
最佳实践是将自动化红队（广度覆盖）
与专业人类红队（深度创新）结合使用。

\textbf{结构化红队框架}
将安全评估组织为系统性的分类体系。
典型的框架按\textbf{危害类别}组织攻击
（暴力内容、非法活动建议、隐私泄露、
偏见歧视、虚假信息传播等），
每个类别包含多个攻击级别
（从简单的直接请求到复杂的多步社会工程）。
这种结构化方法确保了评估的\textbf{完整性}——
不会因为红队成员的个人偏好
而遗漏某些危害类别。

\subsubsection{2026 年的安全评估新进展}

2026 年初，安全评估领域出现了几项
具有理论和实践意义的重要突破。

\textbf{对齐验证不可能定理}（arXiv:2603.08761，2026 年 3 月）
从数学上证明了一个\textbf{根本性的理论极限}：
对于 AI 系统的对齐验证，
\textbf{完备性}（soundness）、
\textbf{通用性}（generality）和
\textbf{可计算性}（tractability）
三者\textbf{无法同时满足}。
这意味着不存在一种既完备、又通用、又高效的方法
来验证 AI 系统是否已经正确对齐。
这一结果类似于 Rice 定理在程序分析中的地位：
它不是说对齐验证不可能，
而是说\textbf{任何实际的验证方法都必须在三者之间做出取舍}。
这对安全评估的实践有深远影响：
我们需要接受「不完美的安全保障」是常态，
并设计能在约束条件下最大化安全性的评估方案。

\textbf{AutoControl Arena}（arXiv:2603.07427）
提出了一种\textbf{自动化}的 frontier AI 风险评估框架，
专门针对\textbf{逻辑幻觉}（logic hallucination）问题，
即模型在推理过程中看似逻辑连贯但实际违反基本推理规则的情况。
与传统的安全评估（主要关注有害内容生成）不同——
AutoControl Arena 关注的是模型在\textbf{高风险决策}场景中
推理错误但自信输出的风险。

\textbf{PersistBench}（arXiv:2602.01146）
首次系统性地评估了 \textbf{Long-Term Memory（LTM）}
带来的安全风险。
随着模型开始支持跨会话的持久记忆，
新的安全威胁浮现：
\begin{itemize}[noitemsep]
  \item \textbf{跨域记忆泄露}：模型将用户 A 的会话信息
        错误地关联到用户 B 的上下文中。
  \item \textbf{记忆诱导的谄媚}（memory-induced sycophancy）：
        模型利用记忆中的用户偏好来调整回答，
        导致越来越迎合用户的错误观点
        而非提供准确信息。
\end{itemize}
PersistBench 为 LTM 安全性建立了首个标准化评估框架。

\textbf{MCP 协议治理}也在 2025 年底迈出了关键一步。
Model Context Protocol（MCP）于 2025 年 12 月
被捐赠给 Linux Foundation 的
\textbf{Agentic AI Foundation}——
成为首个由行业级治理结构管理的 Agent 工具协议。
OpenAI、Google、Microsoft 等主要厂商均已采纳 MCP，
累计下载量超过 9700 万次。
这标志着 Agent 安全从各自为战的防御
走向了标准化的治理框架。

\subsection{对齐税}

安全对齐不是免费的。
让模型更安全通常会\textbf{略微降低}其原始能力。
这被称为\textbf{alignment tax}（对齐税）。

例如，过度的安全训练可能导致模型\textbf{过度拒绝}，
拒绝回答完全合法合理的问题
（如「如何制作鸡尾酒」被错误地归类为有害问题）。
这降低了用户体验。

能力与安全之间存在一个\textbf{Pareto 前沿}，
在给定的技术水平下，不可能同时最大化两者。
更好的对齐技术（如 Constitutional AI、DPO）
可以将这个前沿向外推，
在相同安全水平下保留更多能力，
或在相同能力水平下实现更好的安全性。

安全对齐不是一次性的工作，而是一个持续的过程。
随着模型能力的增强和新攻击方式的出现，
安全机制需要不断迭代更新。

\subsection{Responsible Scaling 与 AI Safety Levels}

2025--2026 年，AI 安全从理论讨论
走向了\textbf{可操作的治理框架}。
Anthropic 率先提出的
\textbf{Responsible Scaling Policy}（RSP）
成为行业参考标准。

RSP 的核心机制是定义一系列
\textbf{AI Safety Levels}（ASL），
类似于生物安全的 BSL 等级：
\begin{itemize}[noitemsep]
  \item \textbf{ASL-1}：无显著风险。
        模型不具备比互联网搜索更多的能力。
        无需特殊安全措施。
  \item \textbf{ASL-2}：当前 frontier 模型的水平。
        模型可能提供一些有害信息
        但不超过互联网上已有的内容。
        需要标准的安全训练和红队测试。
  \item \textbf{ASL-3}：
        模型在生物武器、网络攻击等领域
        提供了\textbf{实质性的能力提升}，
        即帮助攻击者获得
        没有 AI 就无法获得的信息或能力。
        需要严格的访问控制、
        加强的安全评估和部署限制。
  \item \textbf{ASL-4/5}（尚未达到）：
        模型具有自主行动能力、
        可能导致系统性风险的水平。
        需要最高级别的安全措施。
\end{itemize}

RSP 的关键创新在于
它将安全承诺从模糊的意向声明
转化为\textbf{可测量的阈值和可执行的流程}：
模型在发布前必须通过对应 ASL 级别的评估，
评估失败则必须延迟发布并加强安全措施。
OpenAI 随后推出了类似的
\textbf{Preparedness Framework}，
Google DeepMind 发布了
\textbf{Frontier Safety Framework}。
三大实验室的框架在细节上有差异，
但核心逻辑一致：
\textbf{模型能力越强，安全要求越高}。

\subsection{过度拒绝的量化与缓解}

\textbf{过度拒绝}（over-refusal）在 2025--2026 年
成为一个被独立研究的问题。
当模型在安全训练中学到了「拒绝可疑请求」的模式后，
它可能将这个模式\textbf{过度泛化}，
拒绝完全无害的请求，仅仅因为请求的措辞
恰好与有害请求的某些特征相似。

过度拒绝率的量化方法是：
构建一个包含\textbf{表面敏感但实际无害}的请求集合
（如「如何切洋葱不流泪」「炸药奖（诺贝尔奖）的历史」
「如何杀死一个 Linux 进程」），
测量模型错误拒绝的比例。

研究表明，过度拒绝率与安全训练的强度
呈现\textbf{U 形关系}~\cite{cui2024orbench}：
轻度安全训练——拒绝率低但安全性不足；
中度训练——最优平衡；
重度训练——安全性边际收益递减但过度拒绝率飙升。
找到这个「甜蜜点」是后训练团队的核心挑战之一。

\section{评估的实践智慧}

在实际的模型开发过程中，评估团队积累了许多经验教训：

\subsection{持续评估}

\textbf{评估要在训练过程中持续进行}，
而不是等到训练结束后一次性评估。
通常每 1000--5000 步在一组核心基准测试上评估一次，
绘制「能力曲线」，
观察不同能力在训练过程中的变化趋势。
有时会发现某种能力在训练中期突然出现（涌现），
或者某种能力随着训练进行反而退化（灾难性遗忘）。
这些信号可以指导数据配比和训练策略的调整。

\textbf{评估频率与策略分层}是持续评估的关键设计决策。
实践中，团队通常采用多层评估策略：
\textbf{轻量级监控}（每 100--500 步）：跟踪 loss、gradient norm、
learning rate 等训练动态指标，
这些指标计算成本几乎为零，但能快速检测训练不稳定性；
\textbf{核心基准}（每 1000--5000 步）：在 MMLU、GSM8K、
HumanEval 等 5--10 个代表性基准上评估，
每次评估消耗数十到数百 GPU$\cdot$小时；
\textbf{全面评估}（每 10,000--50,000 步或每个训练阶段结束时）：
运行完整的评估套件（通常 30+ 个基准），
包括长上下文、多语言、安全性等维度。
EleutherAI 的 \texttt{lm-evaluation-harness}（截至 2026 年 2 月已更新至 v0.4.11）
是最广泛使用的自动化评估框架，
支持 200+ 个基准的标准化运行。

\textbf{能力曲线分析}（capability curve analysis）
是从持续评估数据中提取训练洞察的核心方法。
通过将不同能力维度的得分绘制为训练步数的函数，
研究者可以识别几类关键模式：
\textbf{线性增长区间}：能力随训练数据量稳定提升，
这通常意味着当前的数据配比和学习率是合理的；
\textbf{涌现拐点}：某项能力在特定步数处突然跃升
（如 chain-of-thought 推理能力往往在 100B+ tokens 后才涌现），
这些拐点的出现时间可以反向验证数据质量；
\textbf{能力退化}：某项能力在训练中后期下降，
通常与数据配比失衡（如代码数据占比下降导致代码能力退化）
或灾难性遗忘有关。
当检测到退化信号时，团队可以动态调整数据混合比例。
这种「数据配比的在线调控」已成为
frontier 模型训练的标准实践。

\textbf{训练失败的早期预警}同样依赖持续评估。
gradient norm 的突然飙升、loss 的非单调波动、
或 perplexity 在验证集上的持续上升，
都是训练即将发散的前兆。
Llama 3 的训练报告~\cite{dubey2024llama3}
详细记录了他们如何通过监控这些信号
在多次训练 spike 中及时干预（回退到上一个健康 checkpoint、
降低学习率或剔除异常数据 shard）。
DeepSeek-V3 更进一步，实现了
基于实时监控的\textbf{自动化故障恢复}：
当系统检测到 loss spike 超过预设阈值时，
自动回退 checkpoint 并隔离可能的故障 GPU 节点，
无需人工干预。这种自动化是实现
“zero irrecoverable loss spikes” 的关键。

\subsection{评估基础设施的工程化}

2025--2026 年，LLM 评估从手动的、临时的流程
演变为\textbf{工程化的自动流水线}。
这种转变的驱动力是双重的：
模型迭代速度的加快
（前沿实验室每周可能产出多个候选 checkpoint）
使得手动评估成为瓶颈；
评估覆盖面的扩大
（从几个核心基准增长到 30+ 个基准套件）
使得自动化成为必需。

\textbf{评估编排系统}是这一工程化的核心组件。
它负责管理评估的完整生命周期：
\begin{enumerate}[noitemsep]
  \item \textbf{触发}：新 checkpoint 保存时自动触发评估，
        或手动触发特定基准的评估。
  \item \textbf{调度}：将评估任务分配到评估集群的空闲 GPU 上，
        按优先级排队（紧急的比较评估优先于例行检查）。
  \item \textbf{执行}：运行标准化的评估脚本，
        确保环境一致性（相同的 prompt 模板、
        相同的 temperature 和 top-p 设置、
        相同的最大生成长度）。
  \item \textbf{聚合}：收集所有基准的结果，
        生成统一格式的评估报告。
  \item \textbf{告警}：当任何能力维度出现
        统计显著的退化时，自动通知团队。
\end{enumerate}

\textbf{评估的可重复性}是一个被严重低估的挑战。
同一模型在「完全相同」的评估条件下
运行两次，分数可能不同。
原因包括：GPU 浮点运算的非确定性、
batch size 对注意力计算的影响、
甚至 Python 字典遍历顺序的随机性。
成熟的评估管道通过
\textbf{固定随机种子}、
\textbf{使用确定性算法}（如 torch 的 deterministic mode）、
以及\textbf{多次运行取平均}
来最大限度地保证可重复性。
但完全的可重复性在大规模评估中
仍然是一个未解决的工程问题。

\textbf{评估成本的优化}是另一个实践重点。
在 70B 模型上运行完整的评估套件
（MMLU 14K 题 + HumanEval 164 题 + GSM8K 8.5K 题 + ...）
可能需要 100+ GPU$\cdot$小时。
为了降低成本，团队采用多种策略：
\textbf{分层评估}（前文已述的轻量/核心/全面三层体系）、
\textbf{子集采样}（在基准的随机子集上评估，
用置信区间估计全集分数，
通常 30\% 的子集就能将估计误差
控制在 $\pm$1\% 以内）、
以及\textbf{早停}（如果前 10\% 的题目表现
明显差于上一版本，提前终止评估并告警）。

\textbf{EleutherAI lm-evaluation-harness}
是评估工程化的开源基础设施。
截至 v0.4.11（2026 年 2 月），
它支持 200+ 个标准化基准，
提供统一的 YAML 配置接口，
并通过插件架构支持自定义评估任务。
它的最大价值在于\textbf{标准化}：
使用 lm-evaluation-harness 的评估结果
在不同团队之间是可比较的，
这在学术界发表论文和
工业界模型选型中都至关重要。
HuggingFace 的 \textbf{Lighteval}
提供了类似的功能但更轻量，
适合快速原型验证。

\subsection{能力维度分析}

\textbf{不同基准测试衡量不同能力维度}。
MMLU 主要测试\textbf{知识记忆}，
即模型在训练数据中是否见过相关知识。
GSM8K 和 MATH 测试\textbf{推理能力}，
即模型能否进行多步逻辑推导。
HumanEval 和 SWE-bench 测试\textbf{代码能力}，
前者是函数级代码生成，后者是系统级 bug 修复。
LMSYS Chatbot Arena 测试\textbf{综合实用性}，
即真实用户在盲评中更偏好哪个模型。

一个模型可能在某些维度上很强，
在另一些维度上很弱。
这就是为什么单一基准测试分数永远不够。
完整的评估需要\textbf{多维度的能力画像}。

\textbf{能力雷达图}（capability radar chart）
是可视化多维度评估结果的标准工具。
将 5--10 个核心能力维度
（知识、推理、代码、多语言、长上下文、指令跟随、安全性等）
作为雷达图的轴线，
可以直观地展示模型的「能力形状」。
这种可视化在模型选型和产品决策中极其有用：
一个「全能型」模型的雷达图接近正多边形，
而一个「特长型」模型则呈明显的不规则形状。
实践中，产品团队往往需要的不是「最强」的模型，
而是在特定维度上最强的模型。
一个医疗问答系统需要知识和安全性的最高分，
而代码能力的得分几乎无关紧要。

\textbf{能力维度之间的相关性与取舍}
也是值得关注的分析方向。
研究表明，某些能力维度之间存在正相关
（如知识记忆与推理能力往往同步提升），
而另一些维度之间可能存在\textbf{张力}。
例如，过度优化安全性（拒绝更多请求）
可能降低模型的实用性（helpful 维度得分下降），
这就是所谓的\textbf{对齐税}
（alignment tax，详见本章安全对齐一节）。
理解这些维度间的结构性关系，
有助于在训练和后训练阶段
做出更有针对性的优化权衡。

\subsection{数据泄露防范}

\textbf{基准测试的「数据泄露」是一个严肃问题}。
如果测试题目出现在训练数据中，
模型可能只是在「背答案」而非真正理解。
防御措施包括：使用 contamination detection
（检查训练数据中是否包含测试题）、
定期更新基准测试、以及使用动态生成的评估问题。

一个实用的检测方法是\textbf{扰动测试}：
对基准测试题做微小的修改（改变数字、替换名称），
如果模型的准确率大幅下降，
说明它可能只是记住了原题的答案，
而没有真正理解解题方法。

\textbf{学术界对污染检测的方法论}
在 2024--2025 年有了显著进展。
Shi 等人~\cite{shi2024detecting}
提出了 Min-K\% Prob 方法：
如果一个文本片段中
概率最低的 K\% token 的平均概率仍然异常高，
则强烈暗示该文本出现在训练数据中。
Golchin \& Surdeanu~\cite{golchin2024time}
则利用 LLM 的自回归特性，
通过让模型「补全」基准测试题的内容
来判断是否存在记忆效应。
这些方法为第三方独立审计提供了工具，
即使无法访问训练数据本身，
也能从模型的行为中推断污染的可能性。

\textbf{从静态基准到动态评估}
是应对数据泄露的根本性解决方案。
LiveBench~\cite{white2024livebench}
通过定期从最新的数学竞赛、新闻和学术论文中
自动生成测试题，
从根本上消除了训练数据包含测试题的可能性。
GPQA~\cite{rein2024gpqa} 通过使用博士级专家
编写的”Google-proof”问题（即使允许搜索引擎也难以找到答案的问题），
提高了「背答案」的难度。
这一趋势在 2025--2026 年进一步加速：
越来越多的评估机构采用
\textbf{一次性使用}（single-use）的评估集，
每个测试集仅用一次，永不公开，
从根本上杜绝了有意或无意的数据泄露。

\subsection{评估金字塔}

\begin{corebox}[评估金字塔：从自动到人工]
\begin{enumerate}[noitemsep]
  \item \textbf{自动化指标}（底层，最便宜）：
        Perplexity、BLEU、ROUGE 等。
        快速、廉价、可重复，但与人类判断相关性有限。
        适合训练过程中的快速检查。
  \item \textbf{基准测试套件}（中层）：
        MMLU、HumanEval、GSM8K 等标准化测试。
        覆盖多个能力维度，有历史可比性。
        适合模型之间的系统性比较。
  \item \textbf{LLM-as-Judge}（中上层）：
        用 GPT-4 等强模型评估其他模型的输出。
        比人类评估便宜得多，与人类判断相关性较高。
        但存在自我偏见（倾向偏好自己风格的输出）。
  \item \textbf{人类评估}（顶层，最贵但最可靠）：
        Chatbot Arena 式的盲评。
        直接衡量用户满意度。
        适合最终的模型质量判断。
  \item \textbf{生产指标}（最终裁判）：
        用户留存率、任务完成率、NPS。
        真实场景中的表现是唯一的最终标准。
\end{enumerate}
\end{corebox}

\subsection{建立针对性评估}

公共基准测试衡量的是通用能力，
但实际应用场景往往有独特的需求。
\textbf{建立针对你自己用例的评估套件}
比追逐公共基准分数更有价值。

例如，如果你在用 LLM 做客服机器人，
你应该建立一个包含真实客户问题的评估集，
用领域专家标注正确答案，
然后衡量模型在\textbf{你的}场景中的表现。
公共基准上多 2 个百分点不如
你的客服评估集上多 5 个百分点有意义。

\textbf{构建针对性评估套件的方法论}
包含几个关键步骤。
首先，\textbf{收集真实用例}，
从生产日志中提取最常见的 100--500 个用户请求，
作为评估集的核心。
其次，\textbf{分层标注}，
将用例按难度和类型分类
（简单事实查询、多步推理、创意生成、敏感话题处理等），
确保评估集覆盖所有关键场景。
第三，\textbf{定义评分标准}，
对每个用例建立明确的评分量表
（如 1--5 分的准确性、完整性、安全性评分），
并用 LLM-as-Judge 实现自动化评分。
Braintrust、DeepEval 等 2025--2026 年涌现的评估平台
极大地降低了构建自定义评估套件的技术门槛，
团队可以在几小时内建立起
针对自身场景的自动化评估管道。

\textbf{评估集的维护}与构建同样重要。
生产场景会随时间变化，
新的用户需求、新的边缘案例、新的安全威胁不断涌现，
评估集必须定期更新以反映这些变化。
一个实用的策略是\textbf{评估集的持续增长}：
每周从生产日志中采样 10--20 个
模型表现不佳的案例加入评估集，
同时淘汰已被所有候选模型完美解决的过时用例。
这种「活的」评估集比静态基准测试
更能准确反映模型在真实场景中的表现。

\subsection{A/B 测试}

在生产环境中，\textbf{A/B 测试是最终的评估方式}。
将两个模型版本随机分配给用户，
比较关键的业务指标，
如用户满意度、任务完成率、对话轮次等。
这是最接近「真相」的评估，
因为它直接衡量了模型在真实场景中的表现。

\textbf{LLM 的 A/B 测试与传统软件的 A/B 测试有本质区别}。
传统 A/B 测试的指标（如点击率、转化率）是确定性的，
同一用户在同一场景下的行为是可重复的。
但 LLM 的输出是\textbf{非确定性}的：
同一个 prompt 在不同请求中可能产生不同的回答，
用户对同一质量回答的评价也可能因上下文而异。
这意味着 LLM A/B 测试需要更大的样本量
才能达到统计显著性。
实践中，通常需要每组至少 1000--5000 次交互，
并使用 bootstrap 或 permutation test
而非简单的 t 检验来评估显著性。

\textbf{在线评估 vs.\ 离线评估的互补性}
是 A/B 测试设计的核心考量。
离线评估（基准测试、LLM-as-Judge）可以在部署前
快速筛选候选模型，但它无法捕捉真实用户行为的复杂性。
例如，一个在离线评估中更「准确」的模型
可能因为回答过于冗长而导致用户流失。
成熟的团队通常建立\textbf{离线-在线相关性模型}：
追踪离线指标与在线指标之间的历史相关性，
用离线指标做快速决策，用在线 A/B 测试做最终验证。
OpenAI 和 Anthropic 的产品团队
已经公开讨论过这种分层评估方法论：
先用内部 Elo 评分（离线）筛选模型版本，
再用生产流量的 A/B 测试（在线）确认提升。

\textbf{生产级 A/B 测试基础设施}需要解决几个工程挑战。
\textbf{流量分配}：用户粒度的随机分流，
确保同一用户在整个实验期间看到同一模型版本
（避免「用户体验不一致」导致的混淆变量）。
\textbf{指标体系}：除了短期指标（单次满意度评分），
还需要跟踪长期指标（用户留存率、周活跃度），
因为模型切换的效果可能需要数周才能完全显现。
\textbf{安全护栏}：A/B 测试必须配合自动化的安全监控，
如果实验组的有害输出率超过阈值，
系统应自动终止实验并回退到基线模型。
\textbf{渐进式发布}（gradual rollout）是常见的风险控制策略：
先对 1\% 的流量测试新模型，
确认安全和质量指标无异常后逐步扩大到 5\%、20\%、100\%。

\subsection{回归测试}

新模型版本在改进某些能力的同时，
可能在其他方面\textbf{退化}，
这被称为\textbf{能力回归}。
例如，一个更强的数学推理模型
可能在创意写作上反而变差了。
建立\textbf{回归测试套件}（一组覆盖所有关键能力的测试用例），
可以在每次模型更新时快速检查是否有能力退化。

\textbf{回归测试的非确定性挑战}
使得 LLM 的回归检测比传统软件更加复杂。
在传统软件中，回归测试是确定性的，
相同输入总是产生相同输出，
测试要么通过要么失败。
但 LLM 的输出具有随机性（即使 temperature=0，
不同硬件和批量大小也可能产生不同输出），
因此单次运行的分数波动不能直接判定为回归。
实践中的解决方案是\textbf{统计回归检测}：
对每个测试用例运行多次（通常 3--10 次），
用置信区间而非单点分数来判断能力变化。
如果新模型在某个能力维度上的 95\% 置信区间
完全低于旧模型的均值，则标记为显著回归。

\textbf{CI/CD 集成}是回归测试工业化的关键一步。
2025--2026 年，越来越多的团队将 LLM 回归测试
嵌入持续集成管道（CI/CD pipeline）：
每次模型更新、prompt 修改或 RAG 配置变更
都自动触发回归测试套件。
DeepEval、Braintrust 等工具
提供了与 GitHub Actions、GitLab CI 的原生集成，
当回归测试检测到能力退化超过预设阈值时，
自动阻止部署并通知开发者。
这种实践将 LLM 质量保障
从「事后发现问题」转变为「事前拦截问题」。

\textbf{LLM-as-Judge 在回归测试中的应用}
正在弥补自动化指标的不足。
传统的回归测试依赖精确匹配或规则化指标
（如 exact match、F1 score），
但这些指标对开放式生成任务（如摘要、对话）无能为力。
2025--2026 年的趋势是使用强模型（如 GPT-4o、Claude）
作为自动化评审者：
给定旧模型和新模型对同一 prompt 的回答，
Judge 模型判断新回答是否在质量上退化。
这种方法的成本远低于人类评估，
且与人类判断的一致性已达到 85--90\%。
但需要注意的是，Judge 模型本身的偏见
（如偏好更长的回答、偏好与自身风格相似的输出）
可能导致系统性的误判，
因此最佳实践是将 LLM-as-Judge
与传统指标和人类抽检结合使用。

\textbf{人类评估仍然是金标准}。
LMSYS Chatbot Arena 的 Elo 排名系统
让真实用户对两个匿名模型的回答做盲评。
截至 2026 年初，已经积累了超过 800 万次人类评估，
是判断模型综合能力最可靠的指标。
它的优势在于：问题来自真实用户需求，
评估标准由用户自定义（不是研究者预设的），
而且动态更新使得「针对性优化」变得困难。

\section{评估的未来方向}
\label{sec:eval-future}

评估方法论的演进速度
正在努力追赶模型能力的增长速度。
几个正在形成的趋势值得关注：

\subsection{从静态基准到交互式评估}

传统基准的根本局限在于它们是\textbf{静态}的：
固定的题目、固定的评估协议、
固定的评分标准。
这使得它们容易被「教」和被「游戏化」。

\textbf{交互式评估}（Interactive Evaluation）
代表了一种根本不同的范式：
不是给模型一道题然后检查答案，
而是让模型在一个\textbf{动态环境}中完成任务，
环境会根据模型的行为做出响应。
WebArena 和 OSWorld 是这种范式的早期实例，
但更激进的版本正在出现。
例如，让模型在一个模拟的办公环境中
完成一天的工作任务
（处理邮件、参加会议、写报告、做决策），
然后评估工作的\textbf{质量和效率}。

交互式评估的优势在于
它更接近 AI 的实际使用场景，
更难被针对性优化。
但它也带来了新的挑战：
环境的设计和实现成本高昂，
评估结果的可重复性更难保证
（环境的随机性增加了方差），
而且评估的计算成本远高于静态基准。

\subsection{自适应评估}

\textbf{自适应评估}（Adaptive Evaluation）
借鉴了计算机自适应测试（CAT）的理念，
根据模型在前几道题上的表现
动态调整后续题目的难度。
如果模型答对了难题，就给更难的；
如果答错了简单题，就降低难度。

这种方法的优势是\textbf{效率极高}，
通常只需要静态评估 20--30\% 的题目量
就能达到相同的评估精度，
因为大量时间不再浪费在
「太简单」或「太难」的题目上。
在 2026 年，一些团队开始实验
用 Item Response Theory（IRT）模型
来同时估计模型的能力参数和题目的难度参数，
使评估更加精准和高效。

\subsection{过程评估而非结果评估}

随着推理模型的普及，
评估的焦点正在从\textbf{最终答案}
转向\textbf{推理过程}。

一个输出了正确答案但推理过程存在逻辑错误的模型，
与一个推理过程严谨但因计算失误导致最终答案有误的模型相比，
哪个「更好」？
在传统的 ORM 评估中，前者得满分后者得零分。
但从实际应用的角度来看，
后者可能更可靠，
因为它的推理能力更强，
计算错误可以通过外部工具（计算器、代码执行）纠正。

\textbf{过程评估}的标准化
是 2026 年评估社区最活跃的研究方向之一。
核心挑战在于：
什么才算「好的推理过程」？
一步到位的直觉跳跃是否不如
冗长但严谨的逐步推理？
不同领域可能有不同的答案。
在数学证明中，严谨性至关重要；
在创业决策中，直觉判断可能更有价值。

\subsection{安全评估的持续化}

安全评估正在从\textbf{发布前的一次性检查}
转变为\textbf{持续的安全监控}。

这种转变的驱动力是：
模型在部署后面对的攻击
与红队在发布前测试的攻击
存在系统性差异。
真实世界的攻击者更有创意、更持久、
更善于利用社区共享的越狱方法。
因此，安全评估不能止步于发布时刻，
而必须\textbf{贯穿模型的整个生命周期}。

\textbf{持续安全监控}的实践包括：
\begin{itemize}[noitemsep]
  \item \textbf{实时攻击检测}：
        监控 API 请求中的异常模式，
        识别可能的越狱尝试。
        当检测到新的攻击模式时，
        自动将其加入对抗训练数据集。
  \item \textbf{定期红队更新}：
        每月运行一次自动化红队评估，
        使用最新的攻击方法库。
        当 ASR（Attack Success Rate）
        超过预设阈值时触发安全补丁。
  \item \textbf{用户报告系统}：
        建立用户反馈渠道，
        收集用户遇到的安全问题
        （有害输出、错误拒绝、隐私泄露等）。
        这些报告是自动化评估无法覆盖的
        \textbf{长尾安全问题}的重要来源。
  \item \textbf{安全版本管理}：
        当发现严重安全漏洞时，
        能够快速回滚到安全版本
        或发布定向修复补丁。
        这要求安全修复能够独立于功能更新进行，
        一个在传统软件中常见但在 LLM 中
        尚未完全标准化的流程。
\end{itemize}

\subsection{评估的民主化}

评估能力曾经是大型 AI 实验室的专利，
只有拥有大量 GPU 和评估团队的组织
才能系统性地评估模型。
但 2025--2026 年的工具进步
正在使评估能力\textbf{民主化}：

\textbf{低成本评估工具}
（如 DeepEval、Ragas、ARES）
使得任何开发者都可以
为自己的 RAG 应用或微调模型
构建自动化评估管道。
这些工具通常在单块 GPU 或纯 CPU 上运行，
使用小型 Judge 模型（如 7B 参数的模型）
替代昂贵的 GPT-4 API，
将评估成本从数千美元降低到数十美元。

\textbf{评估即服务}（Eval-as-a-Service）
平台（如 Braintrust、Weights \& Biases Evaluate）
提供了托管的评估基础设施，
开发者只需上传模型输出和评估标准，
平台自动运行评估并生成报告。
这进一步降低了评估的技术门槛。

评估的民主化意味着
更多的人能够独立验证模型声称的能力。
这对于打击虚假的基准分数声明、
提高 AI 行业的透明度和可信度，
具有深远的意义。

\begin{corebox}[评估的终极问题]
评估方法论的所有进步
最终都指向一个根本性的哲学问题：
\textbf{我们到底在衡量什么？}

基准测试衡量的是特定任务上的准确率。
人类评估衡量的是用户偏好。
安全评估衡量的是攻击成功率。
但这些指标的集合
是否真正捕获了「模型智能」的全貌？

一个在所有基准上得 90\% 的模型
是否比一个在部分基准上得 95\%
但在其他基准上得 80\% 的模型「更智能」？
一个被 800 万用户偏好的模型
是否比一个被 100 位专家偏好的模型「更好」？

这些问题没有标准答案。
它们提醒我们：
\textbf{评估本身就是一种价值判断}。
我们选择衡量什么，
决定了我们认为什么是重要的，
进而影响了模型的发展方向。
评估不仅是技术工具，
更是 AI 发展的方向盘。
\end{corebox}


%% ============================================================

%% BibTeX entries for references cited in this chapter
%% ---------------------------------------------------
%
% @article{hendrycks2021mmlu,
%   title={Measuring Massive Multitask Language Understanding},
%   author={Hendrycks, Dan and Burns, Collin and Basart, Steven and Zou, Andy and Mazeika, Mantas and Song, Dawn and Steinhardt, Jacob},
%   journal={Proceedings of ICLR},
%   year={2021}
% }
%
% @inproceedings{wang2024mmlupro,
%   title={{MMLU-Pro}: A More Robust and Challenging Multi-Task Language Understanding Benchmark},
%   author={Wang, Yubo and Ma, Xueguang and Zhang, Ge and Ni, Yuansheng and Chandra, Abhranil and Guo, Shiguang and Ren, Weiming and Arulraj, Aaran and He, Xuan and Jiang, Ziyan and Li, Tianle and Ku, Max and Wang, Kai and Zhuang, Alex and Fan, Rongqi and Yue, Xiang and Chen, Wenhu},
%   booktitle={Advances in Neural Information Processing Systems (NeurIPS)},
%   year={2024}
% }
%
% @inproceedings{rein2024gpqa,
%   title={{GPQA}: A Graduate-Level Google-Proof Q\&A Benchmark},
%   author={Rein, David and Hou, Betty Li and Stickland, Asa Cooper and Petty, Jackson and Pang, Richard Yuanzhe and Dirani, Julien and Michael, Julian and Bowman, Samuel R.},
%   booktitle={Proceedings of NAACL},
%   year={2024}
% }
%
% @article{glazer2024frontiermath,
%   title={{FrontierMath}: A Benchmark for Evaluating Advanced Mathematical Reasoning in {AI}},
%   author={Glazer, Elliot and Erdil, Ege and Besiroglu, Tamay and Chicharro, Diego and Chen, Evan and Gunning, Alex and Olsson, Caroline Falkman and Denain, Jean-Stanislas and Ho, Anson},
%   journal={Epoch AI},
%   year={2024}
% }
%
% @inproceedings{jimenez2024swebench,
%   title={{SWE-bench}: Can Language Models Resolve Real-World {GitHub} Issues?},
%   author={Jimenez, Carlos E. and Yang, John and Wettig, Alexander and Yao, Shunyu and Pei, Kexin and Press, Ofir and Narasimhan, Karthik},
%   booktitle={Proceedings of ICLR},
%   year={2024}
% }
%
% @inproceedings{white2024livebench,
%   title={{LiveBench}: A Challenging, Contamination-Free {LLM} Benchmark},
%   author={White, Colin and Dooley, Samuel and Roberts, Manley and Pal, Arka and Feuer, Ben and Jain, Siddhartha and Shwartz-Ziv, Ravid and Jain, Neel and Saifullah, Khalid and Narang, Siddartha and others},
%   booktitle={Proceedings of ICLR},
%   year={2025}
% }
%
% @article{li2024stylecontrol,
%   title={Does Style Matter? Disentangling Style and Substance in {Chatbot Arena}},
%   author={Li, Tianle and Angelopoulos, Anastasios and Chiang, Wei-Lin},
%   journal={LMSYS Blog},
%   year={2024}
% }
%
% @article{lightman2024lets,
%   title={Let's Verify Step by Step},
%   author={Lightman, Hunter and Kosaraju, Vineet and Burda, Yura and Edwards, Harri and Baker, Bowen and Lee, Teddy and Leike, Jan and Schulman, John and Sutskever, Ilya and Cobbe, Karl},
%   journal={Proceedings of ICLR},
%   year={2024}
% }
%
% @inproceedings{setlur2025rewarding,
%   title={Rewarding Progress: Scaling Automated Process Verifiers for {LLM} Reasoning},
%   author={Setlur, Amrith and Nagpal, Chirag and Fisch, Adam and Geng, Xinyang and Eisenstein, Jacob and Agarwal, Rishabh and Agarwal, Alekh and Berant, Jonathan and Kumar, Aviral},
%   booktitle={Proceedings of ICLR},
%   year={2025}
% }
%
% @inproceedings{yang2025beyond,
%   title={Beyond the First Error: Process Reward Models for Reflective Mathematical Reasoning},
%   author={Yang, Zhaohui and He, Chenghua and Shi, Xiaowen and Deng, Shihong and Li, Linjing and Yin, Qiyue and Jiang, Daxin},
%   booktitle={Findings of EMNLP},
%   year={2025}
% }
%
% @article{shi2024detecting,
%   title={Detecting Pretraining Data from Large Language Models},
%   author={Shi, Weijia and Ajith, Anirudh and Xia, Mengzhou and Huang, Yangsibo and Liu, Daogao and Blevins, Terra and Chen, Danqi and Zettlemoyer, Luke},
%   journal={Proceedings of ICLR},
%   year={2024}
% }
%
% @article{golchin2024time,
%   title={Time Travel in {LLMs}: Tracing Data Contamination in Large Language Models},
%   author={Golchin, Shahriar and Surdeanu, Mihai},
%   journal={Proceedings of ICLR},
%   year={2024}
% }
%
% @article{ravaut2024contamination,
%   title={A Comprehensive Survey of Contamination Detection Methods in Large Language Models},
%   author={Ravaut, Mathieu and Ding, Bosheng and Jiao, Fangkai and Chen, Hailin and Li, Xingxuan and Zhao, Ruochen and Qin, Chengwei and Xiong, Caiming and Joty, Shafiq},
%   journal={arXiv preprint arXiv:2404.00699},
%   year={2024}
% }
%
% @article{chen2025contamination,
%   title={Benchmarking Large Language Models Under Data Contamination: A Survey from Static to Dynamic Evaluation},
%   author={Chen, Simin and Chen, Yiming and Li, Zexin and Jiang, Yifan and Wan, Zhongwei and He, Yixin and Ran, Dezhi and Gu, Tianle and Li, Haizhou and Xie, Tao and Ray, Baishakhi},
%   booktitle={Proceedings of EMNLP},
%   year={2025}
% }
%
% @article{bai2022constitutional,
%   title={Constitutional {AI}: Harmlessness from {AI} Feedback},
%   author={Bai, Yuntao and Kadavath, Saurav and Kundu, Sandipan and Askell, Amanda and Kernion, Jackson and Jones, Andy and Chen, Anna and Goldie, Anna and Mirhoseini, Azalia and McKinnon, Cameron and others},
%   journal={arXiv preprint arXiv:2212.08073},
%   year={2022}
% }
%
% @misc{openai2026swebenchpro,
%   title={Why {SWE-bench Verified} No Longer Measures Frontier Coding Capabilities},
%   author={{OpenAI}},
%   year={2026},
%   howpublished={\url{https://openai.com/index/why-we-no-longer-evaluate-swe-bench-verified/}}
% }
%
